% concatenated from Kitsios_GenDiagDim.tex, packages.tex
\documentclass[a4paper]{amsart}


\usepackage{amssymb}
\usepackage{mathtools}


%% hyphen %%
\newcommand{\myhyp}{\textnormal{-}\hskip0pt\relax}
\newcommand{\mydhyp}{\textnormal{--}\hskip0pt\relax}


%% Numbering %%

\numberwithin{equation}{section}


%% COLOURS %%

\usepackage[dvipsnames]{xcolor}



\newcommand{\myQED}{\hfill $\scriptstyle \blacksquare$}


%%  OPERATOS AND FUNCTIONS %%

% Caligraphic, math and greek letters %

\newcommand{\N}{\mathbb{N}}
\newcommand{\C}{\mathbb{C}}
\newcommand{\R}{\mathbb{R}}

\newcommand{\Bcal}{\mathcal{B}}
\newcommand{\Ccal}{\mathcal{C}}
\newcommand{\Ecal}{\mathcal{E}}
\newcommand{\Fcal}{\mathcal{F}}
\newcommand{\Gcal}{\mathcal{G}}
\newcommand{\Hcal}{\mathcal{H}}
\newcommand{\Kcal}{\mathcal{K}}
\newcommand{\Lcal}{\mathcal{L}}
\newcommand{\Mcal}{\mathcal{M}}
\newcommand{\Pcal}{\mathcal{P}}
\newcommand{\Ucal}{\mathcal{U}}
\newcommand{\Vcal}{\mathcal{V}}
\newcommand{\Zcal}{\mathcal{Z}}

\newcommand{\e}{\varepsilon}
\newcommand{\ph}{\varphi}
\newcommand{\G}{\Gamma}

\newcommand{\indF}{\mathbf{1}}


% Algebras %

\newcommand{\cont}[1]{C({#1})}
\newcommand{\contO}[1]{C_{0}({#1})}
\newcommand{\contCmp}[1]{C_{\mathrm{c}}({#1})}
\newcommand{\algCenter}[1]{\Zcal({#1})}
\newcommand{\multAlg}[1]{\Mcal({#1})}
\newcommand{\cptPull}[3]{{\mathrm{Cpt}}({#1}, {#2}^{\ast} {#3} )}


% C*-algebras %

\newcommand{\cstar}{C$^\ast$\myhyp{}algebra}
\newcommand{\cstarSub}{C$^\ast$\myhyp{}subalgebra}
\newcommand{\cstarDiag}{C$^\ast$\myhyp{}diagonal}
\newcommand{\cstarPair}{C$^\ast$\myhyp{}pair}
\newcommand{\calg}[1]{\mathrm{C}^\ast({#1})}
\newcommand{\contAlg}[1]{$\cont{#1}$\myhyp{}algebra}
\newcommand{\contOAlg}[1]{$\contO{#1}$\myhyp{}algebra}
\newcommand{\hModule}[1]{${#1}$\myhyp{}Hilbert module}
\newcommand{\bop}[1]{\Bcal({#1})}
\newcommand{\kop}[1]{\Kcal({#1})}


% Coarse geometry %

\newcommand{\Croe}[1]{C^{\ast}_{\mathrm{Roe}}(#1)}
\newcommand{\Cu}[1]{C^{\ast}_{\mathrm{u}}({#1})}
\newcommand{\Cfp}[1]{C^{\ast}_{\mathrm{fp}}({#1})}
\newcommand{\Clc}[1]{C^{\ast}_{\mathrm{lc}}({#1})}
\newcommand{\Ccp}[1]{C^{\ast}_{\mathrm{cp}}({#1})}
\newcommand{\Croesub}[1]{\C_{\mathrm{Roe}}[{#1}]}
\newcommand{\Cusub}[1]{\C_{\mathrm{u}}[{#1}]}
\newcommand{\Cfpsub}[1]{\C_{\mathrm{fp}}[{#1}]}
\newcommand{\Clcsub}[1]{\C_{\mathrm{lc}}[{#1}]}
\newcommand{\Ccpsub}[1]{\C_{\mathrm{cp}}[{#1}]}
\newcommand{\CcpCartan}[1]{\ell^{\infty}_{\mathrm{cp}}({#1})}
\newcommand{\CroeCartan}[1]{\ell^{\infty}_{\mathrm{Roe}}({#1})}
\newcommand{\CfpCartan}[1]{\ell^{\infty}_{\mathrm{fp}}({#1})}

\newcommand{\coaEq}{\overset{\mathrm{ce}}{\sim}}


% General functions %

\newcommand{\id}[1]{\mathrm{Id}_{#1}}

\newcommand{\rAr}{\longrightarrow}
\newcommand{\lAr}{\longlefttarrow}

\newcommand{\prn}[1]{{(#1)}}

\DeclareMathOperator{\solim}{\mathrm{so\myhyp{}}\lim}

\DeclareMathOperator{\img}{Im}
\DeclareMathOperator{\dist}{d}
\newcommand{\diam}[1]{\mathrm{diam} ({#1})}

\newcommand{\norm}[1]{{\left\lVert #1 \right \rVert}}
\newcommand{\abs}[1]{{\left\lvert #1 \right \rvert}}
\newcommand{\normreg}[1]{{\lVert #1  \rVert}}
\newcommand{\normbig}[1]{{\big\lVert #1 \big \rVert}}
\newcommand{\absbig}[1]{{\big\lvert #1 \big \rvert}}
\newcommand{\normBig}[1]{{\Big\lVert #1 \Big \rVert}}
\newcommand{\absBig}[1]{{\Big\lvert #1 \Big \rvert}}
\newcommand{\normBigg}[1]{{\Bigg\lVert #1 \Bigg \rVert}}

\newcommand{\inProd}[2]{{\left\langle {#1}, {#2} \right \rangle}}

\newcommand{\ball}[3]{{B_{#1} (#2,#3)}}

\newcommand{\distR}[2]{\dist_{\mathbb{R}}(#1, #2 )}
\newcommand{\distX}[2]{{\dist_X ( #1, #2 )}}


\newcommand{\supp}[1]{\text{supp}( #1 )}
\newcommand{\prop}[1]{\operatorname{prop} ( {#1} )}

\newcommand{\trans}[1]{{#1}^{\mathrm{T}}}

\newcommand{\quotspace}[2]{\, ^{{#1} } \big/_{#2} \, }


\DeclareMathOperator{\mapTorusOne}{{\mathbb{T}_{Id}}}
\DeclareMathOperator{\mapTorusTwo}{{\mathbb{T}_{Id,Id}}}
\DeclareMathOperator{\mapTorusk}{{\mathbb{T}_{Id,...,Id}}}


\newcommand{\trace}[1]{\operatorname{Tr} \big( #1 \big)}

\newcommand{\rips}[2]{{P_{#1}\big(#2\big)}}

\newcommand{\normalisers}[2]{{\mathcal{N}}_{#2}{\big({#1}\big)}}




% Dimensions %

\newcommand{\dimnuc}[1]{\dim_{\mathrm{nuc}}({#1})}

\newcommand{\ddim}[2]{\dim_{\mathrm{diag}}({#1}\subseteq{#2})}

\newcommand{\ddimn}[3]{\dim_{\mathrm{diag}}({#1}\subseteq{#2}\, ;~{#3})}

\newcommand{\asdim}[1]{\mathrm{asdim}({#1})}

\newcommand{\towdim}[2]{\dim_{\mathrm{tow}}({#1},~{#2})}

% Various %

\newcommand{\unitsimg}[3]{ {#1}\big({#2}\big) \big( e_{#3}\big)}
\newcommand{\unitsimgO}[3]{ {#1}\big({#2}\big) \big( {#3} \big)}

\newcommand{\unMatImg}[3]{{\Ecal_{#3}^{(#1), #2}}}
\newcommand{\gUnMatImg}[3]{{ \mathsf{g}_{#3}^{(#1), #2}}}
\newcommand{\unMatImgO}[1]{{\mathcal{E}_{#1}}}
\newcommand{\gUnMatImgO}[1]{{\mathsf{g}_{#1}}}

\newcommand{\CoProp}[1]{CoP for ${#1}$}


%%% HYPERLINKS %%%
\usepackage[
  colorlinks=true, 
  linkcolor=black,  
  citecolor=Blue,  
  urlcolor=blue, 
  linktocpage=true
]{hyperref}

%%%%% Metadata 
%%%%%%%%%%%%%%%%%%%%%%%%%%%%%%%%%%%%
\hypersetup{
    pdftitle    =   {Generalised Diagonal Dimension and Applications to Large-Scale Geometry},
    pdfauthor   =   {Christos Kitsios},
    pdfsubject  =   {Preprint}
}

%%%%% THEOREM NUMBERING
%%%%%%%%%%%%%%%%%%%%%%%%%%%%%%%%%%%%
    %%% plain style
    \theoremstyle{plain}
    \newtheorem{thm}{Theorem}[section]        
    \newtheorem{propo}[thm]{Proposition}        
    \newtheorem{lem}[thm]{Lemma}          
    \newtheorem{cor}[thm]{Corollary}      
    \newtheorem{qu}[thm]{Question}        
    \newtheorem{conj}[thm]{Conjecture}    
    \newtheorem{prob}[thm]{Problem}       
    \newtheorem{fact}[thm]{Fact}
    \newtheorem{case}[thm]{Case}

    \newtheorem*{thm*}{Theorem}           
    \newtheorem*{propo*}{Proposition}      
    \newtheorem*{lem*}{Lemma}            
    \newtheorem*{cor*}{Corollary}         
    \newtheorem*{qu*}{Question}          
    \newtheorem*{conj*}{Conjecture}       
    \newtheorem*{prob*}{Problem}       
    \newtheorem*{fact*}{Fact}
    \newtheorem*{case*}{Case}

    \newtheoremstyle{indented}
        {\topsep}  % space above
        {\topsep}  % space below
        { \leftskip=\parindent 
            \rightskip=\parindent  } % body font
        {  } % Indent amount
        {\itshape}  % theorem head font
        {.}        % punctuation after name
        { }         % space after head
        { }         % header specification
    
    \theoremstyle{indented}
    \newcounter{claim_counter}[thm]
    \newtheorem{claim}[{claim_counter}]{Claim}    
    \newtheorem*{prfclaim}{Proof of the claim}

    \numberwithin{equation}{section}
    
    %%% definition style
    \theoremstyle{definition}
    \newtheorem{defin}[thm]{Definition}
    \newtheorem{notation}[thm]{Notation}
    \newtheorem{convention}[thm]{Convention}
    
    \newtheorem{defin*}{Definition}
    \newtheorem*{notation*}{Notation} 
    \newtheorem*{convention*}{Convention}
    \newtheorem*{nb}{N.B}
    \newtheorem*{warning}{Warning}

    \theoremstyle{remark}
    \newtheorem{rmk}[thm]{Remark}            
    \newtheorem{example}[thm]{Example}
    \newtheorem{step}{Step}
    \newtheorem*{rmk*}{Remark}   

    %%% letter numbering style
    \theoremstyle{plain}
    \newtheorem{alphthm}{Theorem}        
    \renewcommand*{\thealphthm}{\Alph{alphthm}}
    \newtheorem{alphcor}[alphthm]{Corollary}   
    %\newtheorem{alphdef}[alphthm]{Definition}  
    \newtheorem{alphpropo}[alphthm]{Proposition}
    \newtheorem{alphconj}[alphthm]{Conjecture}
    
    \theoremstyle{definition}
    \newtheorem{alphdef}[alphthm]{Definition}  

%%%%%%%%%%%%%%%%%%%%%%%%%%%%%%%%%%%%




\begin{document}



%%%%% TITLE
%%%%%%%%%%%%%%%%%%%%%%%%%%%%%%%%%%%%
\title[Generalised Diagonal Dimension]{Generalised Diagonal Dimension and applications to large-scale geometry}
\date{April 8, 2026}

%%%%% ABSTRACT
%%%%%%%%%%%%%%%%%%%%%%%%%%%%%%%%%%%%
\begin{abstract}
In this paper, we introduce a generalised diagonal dimension.
We explain why the generalised diagonal dimension extends the notion of diagonal dimension defined by Li, Liao, and Winter, and under which conditions these dimensions coincide. 
We prove permanence properties for the generalised diagonal dimension and compare it with the nuclear dimension. 
We investigate applications of the generalised diagonal dimension in large-scale geometry;
specifically, we show that the generalised diagonal dimension of a noncommutative Cartan subalgebra in the C*-algebra of finite-propagation operators on a uniformly locally finite metric space is equal to the asymptotic dimension of the space. 

\end{abstract}

%%%%% AUTHORS
%%%%%%%%%%%%%%%%%%%%%%%%%%%%%%%%%%%%
\author{Christos Kitsios}
\address{Mathematisches Institut, Georg-August-Universit\"{a}t G\"{o}ttingen, Bunsenstr. 3-5, 37073 G\"{o}ttingen, Germany.}
\email{christos.kitsios@uni-goettingen.de}

\maketitle


Large-scale geometry, also known as coarse geometry, is a framework that studies
``large\myhyp{}scale'' properties of spaces, ignoring their ``local'' structure.
Consequently, we identify spaces that have the same large-scale structure, 
even if they differ locally.
More precisely, we study spaces up to \emph{coarse equivalences}.


For two proper metric spaces $X$ and $Y$ 
we define coarse equivalences 
as follows.
A map
${f \colon X~\rAr~Y}$ is called \emph{controlled} if for each ${r>0}$
there exists ${s>0}$, such that if 
${\dist_X(x,x')\le r}$, 
then 
${\dist_Y(f(x),f(x'))\le s}$.
We say that $X$ and $Y$ are \emph{coarsely equivalent}, and denote ${{X}\coaEq{Y}}$,
if there exist controlled maps ${f \colon X \rAr Y}$ and ${g\colon Y \rAr X}$ and
${C>0}$ such that ${\dist_X({g}\circ{f}(x),x)\le C~\forall~x \in X}$ and ${\dist_Y({f}\circ{g}(y),y)\le C~\forall~y \in Y}$.
Coarse geometry is the study of geometric properties that are invariant under coarse equivalences.
For a thorough treatment of coarse geometry, we refer the reader to  the books~\cite{NowakYu2023LargeScaleGeometry, Roe1996IndexTheoryCoarseGeometryTopology, Roe2003lecturesCrsGeometry}.


Properties that are preserved under coarse equivalences are known as \emph{coarse invariants}. 
The {asymptotic dimension} is an example of a coarse invariant.
It 
was introduced by Gromov for finitely generated groups~\cite{Gromov1993AsymptoticInvariantsInfiniteGroups} and it was extended for coarse spaces in~\cite{Roe2003lecturesCrsGeometry} by Roe.  
We say that a metric space $X$ has \emph{asymptotic dimension} at most $d$ and write ${\asdim{X}\le d}$ if:
for any ${r>0}$, there exists a uniformly bounded covering $\Ucal$ of $X$
with a decomposition 
\[\Ucal = \Ucal^{\prn{0}} \sqcup \hdots \sqcup \Ucal^{\prn{d}}\]
such that ${\dist(U , U') > r}$ for all distinct ${U , U' \in \Ucal^{\prn{i}}}$ and 
each ${i \in \{0,\hdots,d\}}$. 
If $X$ and $Y$ are coarsely equivalent, then it can be shown that ${\asdim{X}=\asdim{Y}}$  
and, therefore, the asymptotic dimension is a coarse invariant. 

\emph{Roe algebras} provide a link between the subjects of large\myhyp{}scale geometry and operator algebras, as they can be interpreted as a \cstar{}ic counterpart to the large\myhyp{}scale geometry of spaces.
Roe first introduced them in order to define higher indices of differential operators on Riemannian manifolds~\cite{Roe1988IndexThm1, Roe1988IndexThm2, Roe1993CoarseCohomologyIndexTheory, Roe1996IndexTheoryCoarseGeometryTopology}.
More precisely,
for a Riemannian manifold $M$ one can define a Hilbert space $L^2(M)$, and,
then
the Roe algebra of $M$, denoted by $\Croe{M}$, 
which is the \cstar{} generated by the locally compact operators of finite propagation in $\bop{L^2(M)}$. 
The K\myhyp{}theory groups of the Roe algebra were used to define higher indices of differential operators. 

Similarly, one can define the Roe algebra of a proper metric space.
To do so, fix an ample geometric module $\Hcal_X$ over X and define the Roe algebra $\Croe{X}$ of $X$ 
to be the \cstarSub{} of $\bop{\Hcal_X}$ that is
generated by
the locally compact operators of finite propagation. 
It should be stressed that different choices of ample geometric modules for $X$ give rise to isomorphic \cstar{}s.
We refer the reader to~\cite{WillettYu2020HigherIndexTheory} for an introduction to Roe algebras of proper metric spaces.

Apart from the Roe algebra, we are also interested in other related operator algebras.
The \cstar{} $\Cfp{X}$ is the \cstarSub{} of $\bop{\Hcal_X}$ generated by the operators of finite propagation and is called \emph{the \cstar{} of operators of finite propagation of $X$}. 
In the literature, finite\myhyp{}propagation operators are also known as \emph{band\myhyp{}dominated operators}~\cite{BragaVignati2023GeldandTypeDuality}.
Moreover, if $X$ is a discrete metric space, then 
we can take the (non\myhyp{}ample) geometric module $\ell^2(X)$ over $X$
and define the \emph{uniform Roe algebra of $X$}, denoted by $\Cu{X}$, to be the
\cstarSub{} of $\bop{\ell^2(X)}$ generated by the operators of finite propagation.
Following the naming convention of~\cite{MartinezVigolo2023RoeViaMods},
we refer to the above \cstar{}s (${\Croe{X},~\Cfp{X}}$ and ${\Cu{X}}$) as \emph{Roe\myhyp{}like algebras}. 
In their work, Mart{\'i}nez and Vigolo ~\cite{MartinezVigolo2023RoeViaMods} provide a unified framework 
to deal with either of the Roe\myhyp{}like algebras via coarse geometric modules.
Roe\myhyp{}like algebras have been studied extensively as operator algebras and for their connections with coarse geometry, see~\cite{BaudierBragaFarahKhukhroVignatiWillett2022UniformRoeAlgRigid, KhukhroLiVigoloZhang2021StructureAsymptoticExpanders, Roe2003lecturesCrsGeometry, Sako2014ProperyAandONLPropery,  WhiteWillett2020CartanUniformRoe, Willett2009NotesPropertyA}.



Following the introduction of Roe algebras, it was 
observed
that coarsely equivalent metric spaces have isomorphic Roe algebras \cite{HigsonRoeYu1993CoarseMayerVietoris, Roe1993CoarseCohomologyIndexTheory}. 
One could further ask whether the converse also holds; this problem is known as \emph{C$^\ast$\myhyp{}rigidity}. 
After the foundational result of {\v{S}}pakula and Willett~\cite{SpakulaWillett2013RigidityRoeAlgebras}, 
many papers have advanced the state of the art
by proving the rigidity in increasingly general settings. 
Notable works include 
\cite{BaudierBragaFarahKhukhroVignatiWillett2022UniformRoeAlgRigid,  BragaFarahVignati2020EmbeddingsUniformRoe, BragaFarahVignati2021UniformRoeCoronas, BragaFarahVignati2022GeneralUniformRoeRigidity}
for the rigidity of uniform Roe algebras 
and~\cite{BragaChungLi2020CoarseBaumConnesandRigidity, LiSpakulaZhang2023MeasuredAsymptoticExpandersRigidity} 
for Roe algebras.
The proof of C$^\ast$\myhyp{}rigidity for bounded geometry spaces 
was completed by Mart{\'i}nez and Vigolo in~\cite{MartinezVigolo2025RigidityBoundedGeometry}:


\begin{alphthm}[{C$^\ast$\myhyp{}rigidity}]
    Let $X$ and $Y$ be proper metric spaces of bounded geometry. 
    Then the following are equivalent:
    \begin{enumerate} \renewcommand{\labelenumi}{(\roman{enumi})}
        \item $X$ and $Y$ are coarsely equivalent,
        \item ${\Croe{X}\cong \Croe{Y}}$,
        \item ${\Cfp{X}\cong \Cfp{Y}}$.
    \end{enumerate}
    Moreover, if $X$ and $Y$ are uniformly locally finite, then the above are also equivalent to the following:
    \begin{enumerate}
        \item[(iv)] $\Cu{X}$ and $\Cu{Y}$ are stably isomorphic, i.e. ${\Cu{X}\otimes \Kcal \cong \Cu{Y} \otimes \Kcal}$.
    \end{enumerate}
\end{alphthm}




Having 
C$^\ast$\myhyp{}rigidity at hand, 
it is natural to ask whether coarse invariants, and particularly the asymptotic dimension, can be detected in Roe\myhyp{}like algebras. 
This was partially addressed by Li, Liao and Winter in~\cite{LiLiaoWinter2023diagDim} as an application of diagonal dimension. 
In their work, they introduced a version of nuclear dimension for diagonal \cstarSub{}s, 
called diagonal dimension.


\begin{alphdef}[{Diagonal Dimension,~\cite[Definition~2.1]{LiLiaoWinter2023diagDim}}] \label{intro:def:diagDim}

        Let ${(D_A \subseteq A)}$ be a \cstarSub{} with $D_A$ abelian. 
    We say ${(D_A \subseteq A)}$ has \emph{diagonal dimension} at most $d$, written as
    \[\ddim{D_A}{A} \le d,\]
    if for any finite subset ${\Fcal \subseteq A}$ and ${\e > 0}$, 
    there exist a finite\myhyp{}dimensional  \linebreak 
    C$^{\ast}$\myhyp{}algebra  
    ${F=F^{(0)} \oplus F^{(1)} \oplus \hdots \oplus F^{(d)}}$ with a maximal abelian $\ast$\myhyp{}subalgebra \linebreak
    ${D_F=D^{(0)} \oplus D^{(1)} \oplus \hdots \oplus D^{(d)}}$ 
    and completely positive maps 
    \[ A \overset{\psi}{\longrightarrow} F \overset{\phi}{\longrightarrow} A\]
    such that
    \begin{enumerate} \renewcommand{\labelenumi}{(\arabic{enumi})}
        \item $\psi$ is contractive,
        \item $\norm{\phi \circ \psi(a) - a } < \e $ for every $a \in \Fcal$,
        \item $\phi^{\prn{i}} \coloneq \phi|_{F^{(i)}}$ is a contractive order zero map, i.e. it preserves orthogonality for each ${i \in  \{0,\hdots,d\}}$,
        \item ${\psi(D_A) \subseteq D_F}$,
        \item ${\phi^{\prn{i}} \big( \normalisers{D^{(i)}}{F^{(i)}}\big) \subseteq \normalisers{D_A}{A}}$.
        
    \end{enumerate}
\end{alphdef}



Let $X$ be a uniformly locally finite metric space. 
Note that the algebra of the multiplication operators 
${\ell^\infty\prn{X}}$ in ${\bop{\ell^2\prn{X}}}$ 
is a maximal abelian $\ast$\myhyp{}subalgebra of the uniform Roe algebra $\Cu{X}$.
Li, Liao and Winter showed that the diagonal dimension of the pair ${(\ell^\infty\prn{X} \subseteq \Cu{X})}$ is equal to the asymptotic dimension of $X$.

\begin{alphthm}[{\cite[Theorem~7.7]{LiLiaoWinter2023diagDim}}] 
    Let $X$ be a uniformly locally finite metric space. Then
    \begin{equation*}
        \ddim{\ell^\infty\prn{X}}{\Cu{X}}=\asdim{X}.
    \end{equation*}
\end{alphthm}

We are interested in computing the asymptotic dimension of a space via its Roe algebra or its \cstar{} of finite\myhyp{}propagation operators.
However, the diagonal dimension and the theory of~\cite{LiLiaoWinter2023diagDim} cannot be applied to these \cstar{}s. 
One important obstruction is that 
a \cstar{} is nuclear when its diagonal dimension is finite~\cite{LiLiaoWinter2023diagDim},
while the Roe algebra and the \cstar{} of finite\myhyp{}propagation operators 
of an ample geometric module are never nuclear.
Another obstruction is that these Roe\myhyp{}like algebras do not have a natural candidate for an abelian $\ast$\myhyp{}subalgebra.
Instead, they have a natural noncommutative Cartan subalgebra~\cite{MartinezVigolo2023RoeViaMods}, which is a noncommutative version of Cartan subalgebras introduced by Exel~\cite{Exel2011noncommCartan}
and further extended by Kwa\'{s}niewski and Meyer \cite{KwasniewskiMeyer2020NoncommutativeCartan}.



Motivated by the problem of detecting the asymptotic dimension in Roe\myhyp{}like algebras,
we introduced a generalisation of the diagonal dimension to overcome the aforementioned obstructions.



\begin{alphdef}[Generalised diagonal dimension] \label{intro:def:genDiagDim}
    Fix a \cstar{} $B$. 
    Let \linebreak
    ${(D_A\subseteq{A})}$ be a nondegenerate \cstarSub{}. 
    We say ${(D_A \subseteq A)}$ has \emph{generalised diagonal dimension with respect to $B$} at most $d$, written as
    \[\ddimn{D_A}{A}{B} \le d,\]
    if for any finite subset 
    ${\Fcal \subseteq {A}}$ and ${\e > 0}$, there exist a finite\myhyp{}
    \linebreak 
    dimensional
    \cstar{} 
    ${F=F^{(0)}\oplus F^{(1)} \oplus \hdots \oplus F^{(d)}}$ 
    with canonical diagonal \linebreak
    ${D_F = D^{(0)} \oplus D^{(1)} \oplus \hdots \oplus D^{(d)}}$ 
    and completely positive maps 
    \[ {A \overset{\psi}{\longrightarrow} F \otimes B \overset{\phi}{\longrightarrow} A}\]
    such that
    {
    \begin{enumerate} \renewcommand{\labelenumi}{(\arabic{enumi})}
        \item $\psi$ is contractive,
        \item ${\normreg{\phi \circ \psi(a) - a } < \e }$ for every $a \in \Fcal$,
        \item ${\phi^{(i)} \coloneq \phi|_{F^{(i)} \otimes B}}$ is a contractive order zero map for each ${i \in \{0,\hdots,d\}}$,
        \item ${\psi(D_A) \subseteq D_F \otimes B}$,
        
        
        \item ${\phi^{(i)} \big( {D^{(i)} \otimes B} \big) \subseteq \normalisers{D_A}{A}}$
        and
        ${\phi^{(i)} ( v \otimes B )  \subseteq \normalisers{D_A}{A}}$
        for each matrix unit ${v \in F^\prn{i}}$ with respect to ${D^{\prn{i}}}$,
        
        \item for each ${i\in \{0,\hdots, d\}}$, if
        $\pi^{(i)}$ is a supporting ${\ast}$\myhyp{}homomorphism for the order zero map $\phi^{(i)}$ and ${\{u_\lambda\}_\lambda}$ is an approximate unit of $B$, then
        \[{\solim_\lambda \pi^{(i)}(v \otimes u_\lambda) \in (D_A)'},\]
        for each minimal projection ${v \in D^{\prn{i}}}$. \label{eq:condition6GenDiagDim}
        
        
    \end{enumerate}
    }
\end{alphdef}


The key idea behind our generalisation of the diagonal dimension 
is to ``enlarge the coefficients'' in the completely positive approximations:
we replace the 
finite\myhyp{}dimensional \cstar{}s with 
finite\myhyp{}dimensional \cstar{}s \emph{over a given \cstar{}},
which is the \cstar{} $B$ in the definition above.
This was inspired by the analogous comparison between the uniform Roe algebra and the \cstar{} of finite\myhyp{}propagation operators on a discrete metric space,
and it allows us to ``escape'' the nuclearity requirement of \cstar{}s with finite diagonal dimension.
It is worthwhile noting that, 
while most of the conditions in Definition~\ref{intro:def:genDiagDim} 
are clear analogues of the conditions in Definition~\ref{intro:def:diagDim}, 
Condition~\eqref{eq:condition6GenDiagDim} is a completely new requirement, which turns out to be important when dealing with infinite dimensional coefficients.
The reason why this condition is useful in our setting is that
when $A$ is a Roe\myhyp{}like algebra on a discrete metric space and $B$ is a unital \cstar{},
Condition~\eqref{eq:condition6GenDiagDim} implies that
if ${u, v\in F^\prn{i}}$ are two distinct minimal projections, then the operators
${\phi^\prn{i}\prn{u\otimes 1_B}}$ and ${\phi^\prn{i}\prn{v\otimes 1_B}}$ are supported on disjoint sets. 
We will use this property
when we compare the generalised diagonal dimension with the asymptotic dimension.

\medskip


It should be stressed that this condition is automatically satisfied for $\C$ coefficients. 
That is,
the generalised diagonal dimension
extends the diagonal dimension of Li, Liao, and Winter.


\begin{alphpropo}
    Let ${(D_A \subseteq A)}$ be a nondegenerate \cstarSub{}, where $D_A$ is abelian.
    Then 
    \begin{equation*}
        \ddim{D_A}{A} = \ddimn{D_A}{A}{\mathbb{C}}.
    \end{equation*}
\end{alphpropo}

\bigskip


Our main result (Theorem~\ref{cor:diagDimFPAlg} below) is that for any uniformly locally finite metric space $X$,  
the generalised diagonal dimension of the noncommutative Cartan subalgebra in the \cstar{} of finite\myhyp{}propagation operators $\Cfp{X}$ is equal to the asymptotic dimension of $X$.

\begin{alphthm} \label{intro:thm:diagDimFPAlg}
     Let $X$ be a uniformly locally finite metric space, let $\Hcal$ be a separable, infinite\myhyp{}dimensional Hilbert space and set ${\Hcal_X \coloneq \ell^2(X, \Hcal)}$. 
     Then
    \begin{equation*}
         \ddimn{\CfpCartan{X}}{\Cfp{\Hcal_X}}{\ell^\infty(\N, \bop{\Hcal})} = \asdim{X}.
    \end{equation*}
\end{alphthm}

\medskip

We expect that a similar result holds for the Roe algebra,
that is, 
\begin{equation*}
         \ddimn{\CroeCartan{X}}{\Croe{\Hcal_X}}{\ell^\infty(\N, \kop{\Hcal})} = \asdim{X},
\end{equation*}
when $X$ is a uniformly locally finite metric space.
However, this has yet to be proved.
One difficulty that arises in proving the above is that the coefficients
${\ell^\infty(\N, \kop{\Hcal})}$ do not form a von Neumann algebra,
while our proof of Theorem~\ref{intro:thm:diagDimFPAlg} uses that ${\ell^\infty(\N, \kop{\Hcal})}$ is a von Neumann algebra.




\bigskip
\noindent\textbf{Structure of the paper.}
The paper is structured as follows. 

In Section~\ref{section:prelim}, we cover some preliminaries on order zero maps, Cartan subalgebras, and Roe-like algebras.

We introduce the generalised diagonal dimension in Section~\ref{section:genDiagDim} and show that it extends the diagonal dimension of Li, Liao, and Winter. We also prove permanence properties of the generalised diagonal dimension and compare it with the nuclear dimension. 

The proof of the main result (Theorem~\ref{intro:thm:diagDimFPAlg}) is presented in Section~\ref{section:diagDimRoeAlg}.



\bigskip
\noindent\textbf{Acknowledgements.}
I would like to thank Federico Vigolo and Ralf Meyer for their guidance throughout this project, and Diego Mart\'{i}nez and Wilhelm Winter for helpful comments.
This work was funded by the RTG 2491 \mydhyp{} \emph{Fourier Analysis and Spectral Theory} of the DFG.


\tableofcontents


%%%%%%%%%%%%%

\section{Preliminaries} \label{section:prelim}

In this paper, all \cstar{}s are assumed to be concrete, that is, they are faithfully represented on some Hilbert space $H$. 


\subsection{Order zero maps}

Winter and Zacharias introduced order zero maps, a class of completely positive maps 
which preserve orthogonality~\cite{WinterZacharias2009orderZero}. 
\begin{defin}[{\cite[Definition~2.3]{WinterZacharias2009orderZero}}]
    Let $A$ and $B$ be \cstar{}s and ${\phi \colon A~\rAr~B}$ be a completely positive map. 
    We say that $\phi$ has \emph{order zero}, if
    \[\phi(a) \phi(a') =0,\]
    for each $0 \le a \in A$ and $0 \le  a' \in A$, satisfying $a a' =0$.
    
\end{defin}

Any $\ast$\myhyp{}homomorphism between \cstar{}s
is easily seen to be an order zero map and 
order zero maps are closely related to $\ast$\myhyp{}homomorphisms. 
Winter and Zacharias in~\cite{WinterZacharias2009orderZero} proved a structure theorem for order zero maps, and defined a $\ast$\myhyp{}homomorphism supporting such maps.

\begin{thm}[Central theorem of order zero maps, {\cite[Theorem~3.3]{WinterZacharias2009orderZero}}] \label{thm:orderZero}
    Let $A$ and $B$ be \cstar{}s and let ${\phi \colon A \rAr B}$ be an order zero map. 
    Set ${C\coloneq \calg{\phi(A)}\subseteq B}$. 
    Then there is a positive element 
    \[h \in \multAlg{C} \cap C',\]
    with ${\normreg{h} = \normreg{\phi}}$,
    where $C'$ is the 
    commutant\footnote{Note that if ${C \subseteq \bop{\Hcal}}$, then ${C' \coloneq \{ T \in \bop{\Hcal} \colon T c = cT  \, \forall \, c \in C\}}$.} 
    of $C$, 
    and a $\ast$\myhyp{}homomorphism
    \begin{align*}
        \pi_\phi \colon A~\rAr~\multAlg{C} \cap \{h\}',
    \end{align*}
    such that
    \begin{equation*}
        \phi(a) = h \pi_\phi (a) ,
    \end{equation*}
    for each $ a \in A$. 
    Moreover, if $A$ is unital, then $h = \phi(1_A)$.
\end{thm}

A $\ast$\myhyp{}homomorphism $\pi_\phi$ defined by the central theorem for the order zero map $\phi$ 
is called a \emph{supporting $\ast$\myhyp{}homomorphism} for $\phi$. 

\begin{rmk} \label{rmk:proofThmOrderZero}
   If $A$ is unital and ${C \subseteq \bop{\Hcal}}$ is nondegenerate, 
  then the supporting $\ast$\myhyp{}homomorphism
    \begin{align*}
        \pi_\phi \colon A~\rAr~\bop{\Hcal}
    \end{align*}
    of Theorem~\ref{thm:orderZero} is
    given by
    \begin{align*}
        \pi_{\phi} (a) = \solim_{n\to \infty} \Big( h + \frac{1}{n} 1_{\Hcal} \Big)^{-1} \phi(a) ,
    \end{align*} for each $a \in A$, where $\solim$ denotes the limit in $\bop{\Hcal}$ with respect to the strong operator topology (SOT).
\end{rmk}

\begin{rmk} \label{rmk:orderZeroSOT}
    Let ${(D_A \subseteq A)}$ and ${(D_B \subseteq B)}$ be nondegenerate \cstarSub{}s
    and
    suppose that
    $A$ is unital and ${1_A \in D_A}$.
    Moreover, suppose that 
    \[\phi \colon A~\rAr~B\] 
    is an order zero map, such that ${\phi(D_A) \subseteq D_B}$. 
    Assume that $B$ acts nondegenerately on the 
    Hilbert space $\Hcal$.
    By Remark~\ref{rmk:proofThmOrderZero},
    we see that
    \begin{equation*}
        \pi_{\phi} (D_A) \subseteq \multAlg{D_B} \cap h' 
        \subseteq \overline{\calg{D_B, 1_{\Hcal}}}^{\mathrm{SOT}},
    \end{equation*}
    where ${\overline{\calg{D_B, 1_{\Hcal}}}^{\mathrm{SOT}}}$ is the closure of ${\calg{D_B, 1_{\Hcal}}}$ in ${\bop{\Hcal}}$ with respect to the strong operator topology (SOT).
\end{rmk}

\smallskip

Using the central theorem for order zero maps, Winter and Zacharias defined a functional calculus for such maps: 

\begin{cor}[{Order zero functional calculus,~\cite[Corollary~3.2]{WinterZacharias2009orderZero}}] \label{cor:orderZeroCalc}
    Let $A$ and $B$ be \cstar{}s, ${\phi \colon A~\rAr~B}$ a contractive order zero map and ${f \in \contO{(0,1]}}$. 
    Set ${C \coloneq \calg{\phi(A)} \subseteq B}$ and $h$, $\pi_{\phi}$ to be as in Theorem~\textup{\ref{thm:orderZero}}. 
    Then the map
    \begin{equation*}
        f(\phi) \colon A~\rAr~C \subseteq B
    \end{equation*}
    given by 
    \begin{equation*}
        f(\phi)(a) = f(h) \pi_{\phi}(a)\; \forall \, a \in A
    \end{equation*}
    is a well\myhyp{}defined completely positive order zero map. 
    
    Moreover, if ${\normreg{f}_{\contO{(0,1]}} \le 1}$, then $f(\phi)$ is contractive.
\end{cor}


\medskip
\subsection{Noncommutative Cartan subalgebras}

\begin{defin}
     ${a \in A}$ is called a 
    \emph{normaliser} of $D$ in $A$
    if ${a D a^\ast, a^\ast D a \subseteq D}$.
    We denote by $\normalisers{D}{A}$ the set of normalisers of $D$ in $A$. 
    The \cstarSub{} $D$ is called \emph{regular} if the set $\normalisers{D}{A}$ generates $A$ as a \cstar{}.
\end{defin}

Renault introduced the notion of Cartan subalgebras for \cstar{}s as maximal abelian, 
regular \cstarSub{}s with a faithful conditional expectation. 

\begin{defin}[{\cite{Renault2008cartanSub}}] \label{def:Cartan}
    Let ${(D \subseteq A)}$ be a \cstarSub{}. We say that $D$ is a \emph{Cartan subalgebra of} $A$ if:
    \begin{enumerate}\renewcommand{\labelenumi}{(\arabic{enumi})}
        \item $D$ is a maximal abelian $\ast$\myhyp{}subalgebra (\emph{masa}) of $A$,
        \item ${(D \subseteq A)}$ is nondegenerate, that is, $D$ contains an approximate unit for $A$,
        \item $D$ is regular,
        \item there exists a \emph{faithful conditional expectation} from $A$ onto $D$,
        that is, a completely positive contractive map ${E \colon A~\rAr~D}$
        with ${E|_D = \mathrm{Id}_D}$ and such that $E$ is injective on the set of positive elements of $A$.
    \end{enumerate}
    Moreover, if $D$ has the \emph{unique extension property} relative to $A$, that is,
    every pure state on $D$ extends uniquely to a pure state on $A$, then $D$ is said to be a \emph{\cstarDiag{}} in $A$.

    For a Cartan subalgebra $D$ in a \cstar{} $A$, we say that the inclusion ${D\subseteq A}$ forms a \emph{Cartan pair}.
\end{defin}

Exel {\cite{Exel2011noncommCartan}} introduced noncommutative Cartan subalgebras  
in separable \cstar{}s by generalising Renault's definition of Cartan subalgebras \cite{Renault2008cartanSub}. 
Building on Exel's work, Kwa\'{s}niewski and Meyer in \cite{KwasniewskiMeyer2020NoncommutativeCartan}
defined noncommutative Cartan subalgebras
without assuming separability.

Following Exel, we introduce virtual commutants in order to define noncommutative Cartan subalgebras.

\begin{defin}[{\cite{Exel2011noncommCartan}}]
    Let ${(D \subseteq A)}$ be a \cstarSub{}. 
    A \emph{virtual commutant} of $D$ in $A$ is a pair
    ${(\phi, J)}$ of a 
    closed
    2\myhyp{}sided ideal ${J \subseteq D}$
    and a linear map ${\phi \colon J~\rAr~A}$ such that
    \begin{equation*}
        \phi(d x) = d \phi(x) 
        \phantom{=} \text{and} \phantom{=}
        \phi(x d) = \phi(x) d,
    \end{equation*}
    for each $x \in J$ and $d \in D$.
\end{defin}

To obtain the definition of noncommutative Cartan subalgebras,
we replace the maximal abelian subalgebra condition from Definition~\ref{def:Cartan}
with a condition on the virtual commutants of the subalgebra.

\begin{defin}[{\cite{Exel2011noncommCartan, KwasniewskiMeyer2020NoncommutativeCartan}}]
    Let ${(D \subseteq A)}$ be a \cstarSub{}. We say that $D$ is a \emph{noncommutative Cartan subalgebra} of $A$ if 
    (2), (3), (4) from Definition~\ref{def:Cartan}, and
    \begin{enumerate}
        \item[(1')] ${\img{\phi} \subseteq D}$ for any virtual commutant ${(\phi,J)}$ of ${D \subseteq A}$
    \end{enumerate}
    hold.
    For a noncommutative Cartan subalgebra $D$ in a \cstar{} $A$, we say that the inclusion ${D\subseteq A}$ forms a \emph{noncommutative Cartan pair}.
\end{defin}


\medskip
\subsection{Roe-like algebras}

In this section, we define Roe algebras 
of discrete metric spaces.
In general, one can define Roe algebras of locally compact, second countable Hausdorff spaces
via geometric modules following
the classical approach of Higson and Roe \cite{HigsonRoe2000AnalyticKHomology},
see also \cite{WillettYu2020HigherIndexTheory} and \cite{MartinezVigolo2023RoeViaMods}.

Throughout this section, ${(X, \dist_X)}$ is a discrete metric space.
   We fix an infinite\myhyp{}dimensional, separable Hilbert space $\Hcal$ and set 
    \[{\Hcal_X \coloneq \ell^2(X, \Hcal)}.\]
    In the sequel, ${\Hcal_X}$ will denote the above Hilbert space for a given discrete metric space $X$, unless stated otherwise.
  For a subset ${U \subseteq X}$, we denote by the characteristic function $\indF_{U}$ of $U$, 
    the 
    projection operator ${\indF_{U} \in \bop{\Hcal_X}}$ given by
    \begin{equation*}
        \indF_{U}(f)(x) 
        =\begin{cases}
            f(x), &x\in U\\
            0, &x\notin U
        \end{cases},
    \end{equation*}
    for each ${f\in \Hcal_X}$.

We can define the support of operators on ${\Hcal_X}$.

\begin{defin}
    Let 
    ${T \in \bop{\Hcal_X}}$. 
    The \emph{support} of $T$, denoted by $\supp{T}$, 
    is the set of all ${(y,x) \in X \times X}$, such that
    \[\indF_V T \indF_U \ne 0\]
    for all open neighbourhoods $U$ of $x$ and $V$ of $y$.
\end{defin}

\begin{defin}
    We define the \emph{propagation} of an operator ${T \in \bop{H_X}}$ as
    \begin{equation*} 
        \prop{T} \coloneq \sup\{ \dist_X(y,x) \colon (y,x) \in \supp{T} \} \in [0, \infty].
    \end{equation*}
    We say that ${T \in \bop{H_X}}$ has \emph{finite propagation}
    if ${\prop{T} < \infty}$.
\end{defin}
    
We now have all the ingredients to define
the \cstar{} of finite\myhyp{}propagation operators.

\begin{defin}
    The \emph{$\ast$-algebra of operators of finite propagation} of $\Hcal_X$ is the $\ast$-algebra
    \begin{equation*}
        \Cfpsub{\Hcal_X} \coloneq \{ T \in \bop{\Hcal_X} \colon \prop{T} < \infty \},
    \end{equation*}
    and its norm closure, denoted by $\Cfp{\Hcal_X}$, is the \emph{\cstar{} of operators of finite propagation of} ${\Hcal_X}$.  
\end{defin}

\begin{rmk}
    Since ${(X, \dist_X)}$ is a discrete metric space,
    each operator ${T \in \bop{\Hcal_X}}$ has an infinite matrix representation given by
    \[T = (T_{x,y})_{x,y \in X} ,\]
    with each $T_{x,y} \in \bop{\Hcal}$. When ${T \in \bop{\Hcal_X}}$ has finite propagation, 
    then
    there exists ${r >0}$, such that
    ${T_{x,y} = 0}$ for each ${x,y \in X}$ with ${\dist_X(x,y) > r}$.
\end{rmk}

\begin{defin} 
    An operator
    ${T \in \bop{\Hcal_X}}$ is \emph{locally compact}
    if 
    \begin{equation*}
        \indF_K T,~T \indF_K \in \kop{\Hcal_X}
    \end{equation*}
    for each compact subset ${K \subseteq X}$.

    The \emph{Roe $\ast$-algebra} of $\Hcal_X$ is the $\ast$\myhyp{}algebra of all locally compact, finite\myhyp{}propagation operators,
    \begin{equation*}
        \Croesub{\Hcal_X} \coloneq \{ T \in \Cfpsub{\Hcal_X} \colon T \text{ is locally compact} \},
    \end{equation*}
    and its norm closure, denoted by ${\Croe{\Hcal_X}}$, is called the \emph{Roe \cstar{}} of $\Hcal_X$, or \emph{Roe algebra} of $\Hcal_X$.
\end{defin}

\smallskip

By taking the Hilbert space $\ell^2(X)$ instead of $\Hcal_X$, we can define the propagation of an operator in ${\bop{\ell^2(X)}}$ similarly.
Then the $\ast$\myhyp{}algebra of operators with finite propagation
\begin{equation*}
        \Cusub{X} \coloneq \{ T \in \bop{\ell^2(X)} ~\colon~\prop{T}<\infty \}
\end{equation*}
is called the \emph{uniform Roe $\ast$\myhyp{}algebra of} $X$, and its norm\myhyp{}closure, which is denoted by 
${\Cu{X}}$, is called the \emph{uniform Roe \cstar{}} of $X$, or \emph{uniform Roe algebra} of $X$.
In the infinite matrix representation picture, the operators 
\[T = (T_{x,y})_{x,y \in X}  \in \bop{\ell^2(X)}\] 
of finite propagation are operators in ${\bop{\ell^2(X)}}$ where ${T_{x,y} \in \C}$, 
and there exists $r >0$ such that
${T_{x,y} = 0}$, for each ${x,y \in X}$ with ${\dist_X(x,y) > r}$.
    
\smallskip

We refer to the Roe algebra, the \cstar{} of operators of finite propagation and the uniform Roe algebra
as \emph{Roe\myhyp{}like algebras},
using the naming convention of~\cite{MartinezVigolo2023RoeViaMods}.


\medskip
\subsubsection{Cartan subalgebras in Roe-like algebras}

Fix a uniformly locally finite metric space ${(X, \dist)}$.
Consider the \cstarSub{} of ${\Cfp{\Hcal_X}}$, given by
\begin{equation}
    \CfpCartan{\Hcal_X} \coloneq \{ T \in \Cfp{\Hcal_X}~\colon~\prop{T} = 0 \}.
\end{equation}
The subalgebra ${\CfpCartan{\Hcal_X}}$ is 
equal to the image of the natural embedding of the algebra 
${\ell^\infty (X, \bop{\Hcal})}$ into ${\Cfp{\Hcal_X}}$,
defined by 
\begin{align*}
    \ell^\infty (X, \bop{\Hcal}) &~\rAr~\Cfp{\Hcal_X}\\
    (T_x)_{x \in X} &~\longmapsto~\, T \colon (\xi_x)_{x \in X}~\longmapsto~(T_x \xi_x)_{x \in X} .
\end{align*}
Similarly, we define the \cstarSub{} of $\Croe{\Hcal_X}$, given by
\begin{equation}
    \CroeCartan{\Hcal_X} \coloneq \{ T \in \Croe{\Hcal_X}~\colon~\prop{T} = 0 \},
\end{equation}
which is equal to ${\ell^\infty (X, \kop{\Hcal})}$ when embedded into ${\Croe{\Hcal_X}}$. 

The \cstarPair{}s
${(\CfpCartan{\Hcal_X}\subseteq\Cfp{\Hcal_X})}$ and ${(\CroeCartan{\Hcal_X}\subseteq\Croe{\Hcal_X})}$ form noncommutative Cartan 
pairs.

\begin{thm}[{\cite[Theorem 6.32]{MartinezVigolo2023RoeViaMods}}]
    Let ${(X, \dist)}$ be a uniformly locally finite metric space.
    Then ${\CfpCartan{\Hcal_X}}$ is a noncommutative Cartan subalgebra of ${\Cfp{\Hcal_X}}$ and
    ${\CroeCartan{\Hcal_X}}$ is a noncommutative Cartan subalgebra of ${\Croe{\Hcal_X}}$.
\end{thm}

The theorem above holds in more generality; however, in this paper, we focus our attention on uniformly locally finite metric spaces.
We refer the reader to~\cite[Section 6.4.]{MartinezVigolo2023RoeViaMods} for more details on noncommutative Cartan subalgebras in Roe\myhyp{}like algebras for coarse spaces of bounded geometry.

\bigskip\bigskip

\section{Generalised diagonal dimension} \label{section:genDiagDim}

Li, Liao, and Winter in \cite{LiLiaoWinter2023diagDim} introduced diagonal dimension, a version of nuclear dimension for diagonal \cstarSub{}s.

\begin{defin}[{Diagonal dimension,~\cite[Definition~2.1.]{LiLiaoWinter2023diagDim}}] \label{def:diagDim}
    Let ${(D_A \subseteq A)}$ be a C$^{\ast}$\myhyp{}subalgebra with $D_A$ abelian. 
    We say ${(D_A \subseteq A)}$ has \emph{diagonal dimension} at most $d$, written as
    \[\ddim{D_A}{A} \le d,\]
    if for any finite subset ${\Fcal \subseteq A}$ and ${\e > 0}$, 
    there exist a finite\myhyp{}dimensional 
    C$^{\ast}$-algebra ${F=F^{(0)} \oplus F^{(1)} \oplus \hdots \oplus F^{(d)}}$ with a maximal abelian $\ast$\myhyp{}subalgebra \linebreak
    ${D_F=D^{(0)} \oplus D^{(1)} \oplus \hdots \oplus D^{(d)}}$ 
    and completely positive maps 
    \[ A \overset{\psi}{\longrightarrow} F \overset{\phi}{\longrightarrow} A\]
    such that
    \begin{enumerate} \renewcommand{\labelenumi}{(\arabic{enumi})}
        \item $\psi$ is contractive,
        \item ${\normreg{\phi \circ \psi(a) - a } < \e} $ for every $a \in \Fcal$,
        \item ${\phi^{\prn{i}} \coloneq \phi|_{F^{(i)}}}$ is a contractive order zero map for each ${i \in  \{0,\hdots,d\}}$,
        \item ${\psi(D_A) \subseteq D_F}$,
        \item ${\phi^{\prn{i}} \big( \normalisers{D^{(i)}}{F^{(i)}}\big) \subseteq \normalisers{D_A}{A}}$.
    \end{enumerate}
\end{defin}

The condition for the normalisers in the definition above can be written 
equivalently in terms of matrix units.
We may replace condition (5) with the following equivalent condition:
\begin{enumerate}
        \item[(5')] $\phi$ maps every matrix unit with respect to $D_F$ into $\normalisers{D_A}{A}$, that is, if ${v\in F}$ such that ${v v^\ast}$ and ${v^\ast v}$ are minimal projections in ${D_F}$, then ${\phi(v) \in \normalisers{D_A}{A}}$.
    \end{enumerate}
The two conditions are equivalent because for finite\myhyp{}dimensional \cstar{}s 
the normalisers of a \cstarDiag{} can be written as a linear combination of pairwise orthogonal matrix units, see e.g.~\cite[Example 2]{Kumjian1986onCDiagonals}.

\begin{rmk} \label{rmk:diagDimVsNucDim}
    By dropping Conditions~(4) and (5) from the definition of the diagonal dimension,
    we obtain the nuclear dimension defined by Winter and Zacharias in \cite{WinterZacharias2010nucDim}.
Therefore, for each nondegenerate \cstarSub{} ${(D_A \subseteq A)}$, we have
\begin{equation*}
    \dimnuc{A} \le \ddim{D_A}{A}.
\end{equation*}
In particular, non\myhyp{}nuclear \cstar{}s always have infinite diagonal dimension.
\end{rmk}



\subsection{Generalised diagonal dimension}

We introduce a generalisation of the diagonal dimension that can take finite values also for non\myhyp{}nuclear \cstar{}s.


\newpage
\begin{defin}[Generalised diagonal dimension] \label{def:genDiagDim}
    Fix a \cstar{} $B$. 
    Let \linebreak
    ${(D_A\subseteq{A})}$ be a nondegenerate \cstarSub{}. 
    We say ${(D_A \subseteq A)}$ has \emph{generalised diagonal dimension with respect to $B$} at most $d$, written as
    \[\ddimn{D_A}{A}{B} \le d,\]
    if for any finite subset 
    ${\Fcal \subseteq {A}}$ and ${\e > 0}$, there exist a finite\myhyp{}
    \linebreak 
    dimensional
    \cstar{} 
    ${F=F^{(0)}\oplus F^{(1)} \oplus \hdots \oplus F^{(d)}}$ 
    with canonical diagonal \linebreak
    ${D_F = D^{(0)} \oplus D^{(1)} \oplus \hdots \oplus D^{(d)}}$ 
    and completely positive maps 
    \[ {A \overset{\psi}{\longrightarrow} F \otimes B \overset{\phi}{\longrightarrow} A}\]
    such that
    {
    \begin{enumerate} \renewcommand{\labelenumi}{(\arabic{enumi})}
        \item $\psi$ is contractive,
        \item ${\normreg{\phi \circ \psi(a) - a } < \e }$ for every $a \in \Fcal$,
        \item ${\phi^{(i)} \coloneq \phi|_{F^{(i)} \otimes B}}$ is a contractive order zero map for each ${i \in \{0,\hdots,d\}}$,
        \item ${\psi(D_A) \subseteq D_F \otimes B}$,
        
        
        \item ${\phi^{(i)} \big( {D^{(i)} \otimes B} \big) \subseteq \normalisers{D_A}{A}}$
        and
        ${\phi^{(i)} ( v \otimes B )  \subseteq \normalisers{D_A}{A}}$
        for each matrix unit ${v \in F^\prn{i}}$ with respect to ${D^{\prn{i}}}$,
        
        \item for each ${i\in \{0,\hdots, d\}}$, if
        $\pi^{(i)}$ is a supporting ${\ast}$\myhyp{}homomorphism for the order zero map $\phi^{(i)}$ and ${\{u_\lambda\}_\lambda}$ is an approximate unit of $B$, then
        \[{\solim_\lambda \pi^{(i)}(v \otimes u_\lambda) \in (D_A)'},\]
        for each minimal projection ${v \in D^{\prn{i}}}$.
    \end{enumerate}
    }
\end{defin}

We say that ${(F,~D_F,~\psi,~\phi)}$ is a \emph{completely positive (c.p.) approximation witnessing 
${\ddimn{D_A}{A}{B} \le d}$ for $\Fcal$ within $\e$} if the tuple ${(F,~D_F,~\psi,~\phi)}$ satisfies the conditions in Definition~\ref{def:genDiagDim} with respect to the finite subset ${\Fcal \subseteq A}$ and ${\e>0}$. 

\smallskip

It can be proved that the completely positive maps in Definition~\ref{def:genDiagDim} can be chosen such that ${\phi \circ \psi}$ is contractive.
For later use, we record below a simple lemma. Its proof is explained in \cite[Remark~2.2.(ii)]{LiLiaoWinter2023diagDim} and \cite[Remark~2.2.(iv)]{WinterZacharias2010nucDim}.

\begin{lem} \label{lem:cpPhioPsiContractive}
    Suppose that $A$ is unital and ${\ddimn{D_A}{A}{B} = d < \infty}$. Then we can choose a completely positive approximation 
    such that ${\phi \circ \psi}$ is contractive.
\end{lem}

\begin{rmk}
      The lemma above can be extended for non\myhyp{}unital \cstar{}s by using an approximate unit, see~\cite[Remark 2.2.(iv)]{WinterZacharias2010nucDim} for similar techniques.
\end{rmk}

\medskip
\subsubsection{Commutation property}

We introduce the following property. 

\begin{defin}[Commutation property CoP] \label{def:PropP}
    Let $B$ be a unital C$^{\ast}$-algebra, $\Hcal$ be a Hilbert space and ${\pi \colon M_r(B)~\rAr~\bop{\Hcal}}$ be a ${\ast}$\myhyp{}homomorphism with ${r \in \mathbb{N}}$. 
    Suppose that ${D \subseteq \bop{\Hcal}}$ is a \cstar{}.
    
    Then we say that $\pi$ has the \emph{commutation property for $D$} (\emph{\CoProp{D}}) 
    if for each minimal projection ${v \in D_r(\C)}$ 
    we have
    \[{d \pi(v\otimes 1_B)  =\pi(v\otimes 1_B)  d},\]
    for each ${d \in D}$.
    
    When $B$ is not unital, we fix an (increasing) approximate unit ${\{u_\lambda\}_{\lambda}}$ of $B$.
    We say that $\pi$ has \emph{\CoProp{D}} if
    for each minimal projection ${v \in D_r(\C)}$ and each ${d \in D}$ it holds
    \[d  \big(\solim_\lambda \pi(v \otimes u_\lambda) \big) = \big( \solim_\lambda \pi(v \otimes u_\lambda)  \big) d,\]
    where $\solim$ denotes the limit in the strong operator topology (SOT).

    We say that ${\pi \colon M_{r_1}(B)  \oplus \hdots \oplus  M_{r_m}(B)~\rAr~\bop{\Hcal}}$ has \emph{\CoProp{D}} if for each ${i\in \{1,\hdots,m\}}$ the restriction ${\pi|_{M_{r_i}(B)}}$ has \CoProp{D}.
\end{defin}

Using the above property we can rephrase Condition~(6) of Definition~\ref{def:genDiagDim} with the following condition: 
for each ${i\in \{0,\hdots, d\}}$, if
$\pi^{(i)}$ is a supporting ${\ast}$\myhyp{}homomorphism 
for the order zero map $\phi^{(i)}$, 
then $\pi^{(i)}$ has \CoProp{D_A}.

\begin{rmk}
    This commutation property is of interest
    because in the sequel, we will define order zero maps whose supporting $\ast$\myhyp{}homomorphisms 
    map the diagonal matrix units to operators supported in pairwise disjoint sets.
    Maps into \cstar{}s with CoP for their noncommutative Cartan subalgebra have this property.
\end{rmk}

\begin{example}
    Let ${\Hcal}$ be an infinite\myhyp{}dimensional, separable Hilbert space and define
    the $\ast$\myhyp{}homomorphism ${\pi}$ given by
    \[\begin{array}{ccccc}
        {\pi} &\colon & M_2(\bop{\Hcal}) &~\longrightarrow~& \Cfp{\ell^2(\N, \Hcal)}\\
        \\
         & & \begin{pmatrix}
            T_{1,1}  & T_{1,2}  \\
            T_{2,1}  & T_{2,2}
        \end{pmatrix}
        &~\longmapsto~& \begin{pmatrix}
            T_{1,1} & T_{1,2} & 0       & 0  & \hdots \\
            T_{2,1} & T_{2,2} & 0       & 0  \\
            0       & 0       & T_{1,1} & T_{1,2}  \\
            0       & 0       & T_{2,1}  & T_{2,2} \\
            \vdots & &  & & \ddots 
        \end{pmatrix} .
    \end{array}\]
    Then ${\pi}$ has \CoProp{\CfpCartan{\ell^2(\N, \Hcal)}}
    and,
    moreover, maps diagonal matrix units to operators with pairwise disjoint supports.
    
    Conversely,
    we can define a
     $\ast$\myhyp{}homomorphism 
    \[\begin{array}{ccccc}
        \widetilde{\pi} & \colon & M_2(\bop{\Hcal}) 
        &~\rAr~& 
        \Cfp{\ell^2(\N, \Hcal)}
    \end{array}\]
    that does not have either of these properties:
    Let ${p, q \in \bop{\Hcal} \backslash \{1_{\Hcal} , 0\}}$ be two infinite\myhyp{}rank projections such that ${q = 1_{\Hcal} - p }$. 
    We can define isometries ${v_1, v_2 \in \bop{\Hcal}}$, satisfying ${v_1 \Hcal = p\Hcal}$ and ${v_2 \Hcal = q\Hcal}$.
    Set
    \[\begin{array}{ccccc}
        v \coloneq v_1 \oplus v_2& \colon & \Hcal  \oplus  \Hcal & \rAr & \Hcal\\
        & & (h_1, h_2 ) & \longmapsto &  v_1 h_1 + v_2 h_2 
    \end{array}.\] 
    Let ${\widetilde{\pi}}$
    be the $\ast$\myhyp{}homomorphism given by
    \[\begin{array}{ccccc}
       \widetilde{\pi} & \colon & M_2(\bop{\Hcal}) 
        &~\rAr~& 
        \Cfp{\ell^2(\N, \Hcal)}\\
        \\
        &  & T \phantom{\Hcal))}
        &~\longmapsto~&
        \begin{pmatrix}
            v T v^\ast  & 0             & 0 & \hdots\\
            0           & v T v^\ast    & 0 &  \\
            0           & 0             & v T v^\ast     &  \\
            \vdots       &     &  & \ddots 
        \end{pmatrix} .
    \end{array}\]
    It is straightforward to see that
    \[\widetilde{\pi} \begin{pmatrix}
        1_{\Hcal} & 0 \\
        0 & 0 
    \end{pmatrix}  \notin \big(\CfpCartan{\ell^2(\N, \Hcal}\big)'\]
    and, hence,
    ${\widetilde{\pi}}$ does not have 
    \CoProp{\CfpCartan{\ell^2(\N, \Hcal)}}.
    Moreover, 
    we have that the operators
   \[\widetilde{\pi} \begin{pmatrix}
        1_{\Hcal} & 0 \\
        0 & 0 
    \end{pmatrix} = \begin{pmatrix}
            p  & 0             & 0 & \hdots\\
            0           & p    & 0 &  \\
            0           & 0             & p    &  \\
            \vdots       &     &  & \ddots 
        \end{pmatrix} 
        \phantom{==} \text{ and } \phantom{==} 
        \widetilde{\pi} \begin{pmatrix}
        0 & 0 \\
        0 & 1_{\Hcal}
    \end{pmatrix} = \begin{pmatrix}
            q  & 0             & 0 & \hdots\\
            0           & q    & 0 &  \\
            0           & 0             & q    &  \\
            \vdots       &     &  & \ddots 
        \end{pmatrix}\]
        are supported on the same set.
        Hence,
        ${\widetilde{\pi}}$ maps the diagonal matrix units to operators supported on the same set.
\end{example}


\medskip
\subsubsection{Comparison with the diagonal dimension}

The generalised diagonal dimension extends the diagonal dimension of Li, Liao, and Winter (Definition~\ref{def:diagDim}).
More precisely, if $D_A$ is an abelian \cstarSub{} in a \cstar{} $A$, then
the generalised diagonal dimension with respect to $\C$ is equal to the diagonal dimension.

\begin{propo} \label{prop:genDiagDimExtendsDiagDim}
    Let ${(D_A \subseteq A)}$ be a nondegenerate \cstarSub{}, where $D_A$ is abelian.
    Then 
    \begin{equation*}
        \ddim{D_A}{A} = \ddimn{D_A}{A}{\mathbb{C}}.
    \end{equation*}
\end{propo}

\begin{proof}
    The inequality $\le$ is automatic. We elaborate on the converse
    inequality.

    From the discussion after Definition~\ref{def:diagDim}, we see that Condition~(5) 
    from Definition~\ref{def:genDiagDim} for $B=\C$
    is equivalent to Condition~(5) in Definition~\ref{def:diagDim}. 
    Therefore, it is only left to show that Condition (6) is automatic when $D_A$ is abelian and $B=\C$.
    
    If ${\phi \colon M_r(\C) \rAr A}$ and ${\phi(D_r(\C)) \subseteq D_A}$, 
    then by the central theorem of order zero maps (Theorem~\ref{thm:orderZero}) 
    we can define a supporting  $\ast$\myhyp{}homomorphism 
    \[\pi \colon M_r(\C)~\rAr~\bop{H}\]
    such that ${\phi\prn{\cdot} = h \pi \prn{\cdot}}$
    and
    ${h \in \big(\img{\pi}\big)'}$. 
    Remark~\ref{rmk:orderZeroSOT} implies
    \[\pi(D_r(\C)) \subseteq \overline{\calg{D_A, \mathbf{1}_{\bop{H}}}}^{\mathrm{SOT}}, \]
    where:
    \begin{itemize}
        \item ${\calg{D_A, \mathbf{1}_{\bop{H}}} }$ is the \cstar{} generated  by $D_A$ and ${\mathbf{1}_{\bop{H}}}$,
        \item ${\overline{\calg{D_A, \mathbf{1}_{\bop{H}}}}^{\mathrm{SOT}}}$ 
        denotes the closure of ${\calg{D_A, \mathbf{1}_{\bop{H}}}}$ in $\bop{H}$ with respect to the strong operator topology (SOT).
    \end{itemize}

    By the von Neumann double commutant theorem, we have
    \[ \overline{\calg{D_A, \mathbf{1}_{\bop{H}}}}^{\mathrm{SOT}} =  \big({\calg{D_A, \mathbf{1}_{\bop{H}}}}\big)'' = {D_A}'',\] where ${D_A}''$ is the double commutant of ${D_A}$.
    Combining the above we obtain
     \[{\pi(D_r(\mathbb{C})) \subseteq D_A{''}}.\] 
    
    If ${u \in D_A}$ unitary, we have that ${u \in {D_A}'}$, since $D_A$ is abelian. 
    Hence, for each ${k \in \{1,\hdots,r\}}$ it holds
    \[  u \pi(e_{k,k}) u^{\ast} =   \pi(e_{k,k}) u u^{\ast} = \pi(e_{k,k}), \]
    where we used that ${\pi(e_{k,k}) \in {D_A}''}$ and ${u \in {D_A}'}$. Therefore, $\pi$ has \CoProp{D_A}, in other words,
    Condition~(6) of Definition~\ref{def:genDiagDim} is satisfied. \qedhere  
\end{proof}

\medskip
\subsection{Permanence properties and nuclear dimension}

The following permanence properties for the generalised diagonal dimension hold.

\begin{propo} \label{prop:permProp}
    Let ${(D_A \subseteq A)}$ and ${(D_C \subseteq C)}$ be nondegenerate \cstarSub{}s.
    Then we have the following properties.
   { \begin{enumerate} \renewcommand{\labelenumi}{(\roman{enumi})}
        \item Direct sums:
    \begin{align*}
        &\ddimn{D_A\oplus D_C}{A \oplus C }{B} \\
       &~\le \max \{ 
       \ddimn{D_A}{A}{B}, 
       \ddimn{D_C}{C}{B} \} .
    \end{align*}
        \item Tensor products:
        \begin{align*}
       &\ddimn{D_A\otimes D_C}{A \otimes C }{B} +1 \\
       &~\le \big( \ddimn{D_A}{A}{\C}+1\big) \big(\ddimn{D_C}{C}{B} +1 \big).
    \end{align*}
    In particular, for each $n \in \N$
    \begin{align*}
        \ddimn{D_n(D_A)}{M_n(A)}{B} \le \ddimn{D_A}{A}{B}.
    \end{align*}
    \end{enumerate}
}
\end{propo}

These permanence properties can be proved by making the natural choices
for the completely positive maps $\phi$ and $\psi$
(direct sums of approximations for (i), and tensor products for (ii)). 
A detailed proof can be found in \cite{Kitsios2025PhDThesis}. 
For similar ideas we refer the reader to the permanence properties
of the nuclear dimension~\cite{WinterZacharias2010nucDim} and 
diagonal dimension~\cite{LiLiaoWinter2023diagDim}.

\smallskip

Using Remark~\ref{rmk:diagDimVsNucDim},
we observe that 
for each nondegenerate \cstarSub{} 
\linebreak
${(D_A \subseteq A)}$, we have
\begin{equation*}
    \dimnuc{A} \le \ddim{D_A}{A} \le \ddimn{D_A}{A}{\C}.
\end{equation*}
Moreover, when $B$ is chosen to be a \cstar{} of finite nuclear dimension,
one can obtain a bound for the nuclear dimension of $A$ 
involving the generalised diagonal dimension.
 
\begin{thm}
    Let ${(D_A \subseteq A)}$ be a nondegenerate \cstarSub{} and let $B$ be a \cstar{} with finite nuclear dimension. 
    If ${\ddimn{D_A}{A}{B}< \infty}$, then $A$ is a nuclear \cstar{} with finite nuclear dimension given by
    \begin{equation*}
        \dimnuc{A} + 1 \le \big(\dimnuc{B}+1\big)\big(\ddimn{D_A}{A}{B}+1\big).
    \end{equation*}
\end{thm}

One can prove the above theorem by defining
completely positive maps for the nuclear dimension of $A$
as direct sums of completely positive approximations for ${\dimnuc{B}}$ and ${\ddimn{D_A}{A}{B}}$, 
see \cite{Kitsios2025PhDThesis} for details.


\bigskip
\section{Diagonal dimension of Roe\myhyp{}like algebras} \label{section:diagDimRoeAlg}

In this section, we 
compute the diagonal dimension of the inclusion of the noncommutative Cartan subalgebra in the \cstar{} of finite\myhyp{}propagation operators on a uniformly locally finite metric space.

Before computing the diagonal dimension we define the asymptotic dimension for metric spaces. 

\begin{defin} \label{def:asDim}
    Let ${(X, \dist)}$ be a metric space.
    Then
$X$ has \emph{asymptotic dimension ${d \ge 0}$}, if $d$ is the least integer,
with the following property:
for any ${r>0}$, there exists a covering ${\Ucal = \{U_j\}_{j \in J}}$ of $X$,
with a decomposition 
\[\Ucal = \Ucal^{\prn{0}} \sqcup \hdots \sqcup \Ucal^{\prn{d}},\]
such that
$\Ucal$ is uniformly bounded, i.e.
\[\sup_{U \in \Ucal} \diam{U}< \infty, \]
and for each ${i \in \{0,\hdots,d\}}$ the family
${\Ucal^{\prn{i}}}$ is $r$\myhyp{}separated, i.e.
${\dist(U_{j} , U_{j'}) > r}$ for each ${U_{j} , U_{j'} \in \Ucal^{\prn{i}}}$ with ${j \ne j'}$. 
\end{defin}

The asymptotic dimension was introduced by Gromov~\cite{Gromov1993AsymptoticInvariantsInfiniteGroups}
for finitely generated groups and 
was later extended by Roe~\cite{Roe2003lecturesCrsGeometry} for a larger class of spaces, including coarse spaces.

\smallskip

We prove that the generalised diagonal dimension in the \cstar{} of finite\myhyp{}propagation operators is equal to the asymptotic dimension of the space.

\begin{thm} \label{cor:diagDimFPAlg}
     Let $X$ be a uniformly locally finite metric space. 
     Then
    \begin{equation*}
         \ddimn{\CfpCartan{X}}{\Cfp{\Hcal_X}}{\ell^\infty(\N, \bop{\Hcal})} = \asdim{X}.
    \end{equation*}
\end{thm}

We split the theorem in two parts, a theorem for an upper bound for the generalised diagonal dimension:
\begin{equation*}
         \ddimn{\CfpCartan{X}}{\Cfp{\Hcal_X}}{\ell^\infty(\N, \bop{\Hcal})} \le \asdim{X},
\end{equation*}
and a theorem for a lower bound:
\begin{equation*}
         \ddimn{\CfpCartan{X}}{\Cfp{\Hcal_X}}{\ell^\infty(\N, \bop{\Hcal})} \ge \asdim{X}.
\end{equation*}
The proof of these inequalities can be found in the following subsections. By combining the upper and lower bound, we deduce Theorem~\ref{cor:diagDimFPAlg}.


\medskip
\subsection{Upper bound}

\begin{thm} \label{thm:upperBoundFPAlg}
    Let $X$ be a uniformly locally finite metric space.
    Then
    \begin{align*}
         \ddimn{\CfpCartan{X}}{\Cfp{\Hcal_X}}{\ell^\infty(\N, \bop{\Hcal})} \le \asdim{X}.
    \end{align*}
\end{thm}

This theorem is obtained by extending the proof of the corresponding inequality in \cite[Theorem~7.7]{LiLiaoWinter2023diagDim}. 
It is worth noting that \cite[Theorem~7.7]{LiLiaoWinter2023diagDim} in turn follows from the proof of \cite[Theorem~8.5]{WinterZacharias2010nucDim}, which states that
\begin{equation*}
    \dimnuc{\Cu{X}} \le \asdim{X}.
\end{equation*}
We only sketch the proof of the theorem here and refer the reader to a detailed proof found in \cite[Theorem~5.2]{Kitsios2025PhDThesis}.

\begin{proof}[Sketch of proof]
    Suppose that $d \coloneq \asdim{X} < \infty$.
    Let $\Fcal \subseteq \Cfp{\Hcal_X}$ be a finite set and $\e>0$.

    Following \cite[Theorem~7.7]{LiLiaoWinter2023diagDim} we may fix ${r \in \N}$ large enough and choose
     ${(d+1)}$ uniformly bounded 
     $3r$\myhyp{}separated families ${\mathcal{U}^{(0)},\hdots,\mathcal{U}^{(d)}}$
     covering $X$, i.e.
     \[X= \bigcup_{i=0}^d \mathcal{U}^{(i)}.\]

     Let
     \begin{equation*}
         B^{(i)} \coloneq \prod_{U \in \mathcal{U}^{(i)}} \indF_{\ball{}{U}{r}} \,  \Cfp{\Hcal_X} \, \indF_{\ball{}{U}{r}},
     \end{equation*}
     and
     \begin{equation*}
         D^{(i)} \coloneq  \prod_{U \in \mathcal{U}^{(i)}} \indF_{\ball{}{U}{r}} \,  D_A \, \indF_{\ball{}{U}{r}},
     \end{equation*}
    and set
      \begin{align*}
        f_i \coloneq \frac{1}{r} \sum_{U \in \mathcal{U}^{(i)}} \sum_{m=1}^{r} \indF_{\ball{}{U}{m}},
    &&
        f \coloneq \sum_{i=0}^d f_i,
   && \text{and}
   &&
        h_i \coloneq \big( f_i f^{-1} \big)^{1/2}, 
    \end{align*}
    for each ${i \in \{0,\hdots,d\}}$.

    We define completely positive maps
    \begin{equation} \label{eq:upperBoundAsdimApprox1}
        A 
        \xlongrightarrow{  \Psi_r}  
       B^{(0)} \oplus B^{(1)} \oplus \hdots \oplus B^{(d)} 
        \xlongrightarrow{ \Phi_r } 
         A
    \end{equation}
    given by
    \[\begin{array}{cccccccccc}
         \Psi_r \colon 
         & A 
         &~\rAr~
         & B^{(0)} 
         & \oplus 
         & B^{(1)} 
         & \oplus & \hdots & \oplus 
         & B^{(d)}\\
         & T 
         &~\longmapsto~
         & (h_0 T^{\prn{0}} h_0
         &  ,
         & h_1 T^{\prn{1}} h_1 
         & ,
         & \hdots & ,
         & h_d T^{\prn{d}} h_d )
     \end{array}\]
     and
         \[ \begin{array}{cccccccccc}
         \Phi_r  \colon & B^{(0)} & \oplus & B^{(1)} & \oplus & \hdots & \oplus & B^{(d)} &~\rAr~& A \\
        & (a_0 & , & a_1 & , & \hdots & , & a_d)
        &~\longmapsto~& a_0 + a_1 + \hdots + a_d    
     \end{array}.
\]
    We see that ${\Psi_r}$ and ${\Phi_r}$ satisfy the following:
    \begin{enumerate} \renewcommand{\labelenumi}{(\alph{enumi})}
        \item \label{eq:upperBoundCPAlgPropa} $\Psi_r$ is a completely positive contractive map,
        \item \label{eq:upperBoundCPAlgPropb} ${\normreg{\Phi_r \circ \Psi_r(T) - T} < \e}$ for each ${T \in \Fcal}$,
        \item \label{eq:upperBoundCPAlgPropc}$\Phi_r|_{B^{(i)}}$ is a contractive order zero map for each ${i \in \{0,\hdots,d\}}$,
        \item \label{eq:upperBoundCPAlgPropd} ${\Psi_r(D_A) \subseteq \bigoplus_{i=0}^d D^{\prn{i}}}$,
        \item \label{eq:upperBoundCPAlgPrope} ${\Phi_r|_{B^{(i)}} \big( \normalisers{D^{(i)}}{B^{(i)} } \big) \; \subseteq  \; \normalisers{D_A }{ A}}$ for each ${i \in \{0,\hdots,d\}}$.
    \end{enumerate}
    It should be stressed that the property (b) above is only true when $r$ is large enough.
    
    The maps ${\Psi_r}$ and ${\Phi_r}$ induce completely positive maps
    \begin{equation*}
        \Cfp{\Hcal_X} 
        \xlongrightarrow{ \psi}  
       \bigoplus_{i=0}^d F^{(i)} \otimes \ell^\infty(\N, \bop{\Hcal})
        \xlongrightarrow{\phi } 
         \Cfp{\Hcal_X},
    \end{equation*}
    where ${F^{(i)}}$ is a finite\myhyp{}dimensional \cstar{} for each ${i \in \{0, \hdots, d\}}$.

    Properties \hyperref[eq:upperBoundCPAlgPropa]{(a)},
    \hyperref[eq:upperBoundCPAlgPropb]{(b)},
    \hyperref[eq:upperBoundCPAlgPropc]{(c)},
    \hyperref[eq:upperBoundCPAlgPropd]{(d)}
    and \hyperref[eq:upperBoundCPAlgPrope]{(e)} imply
    Conditions (1), (2), (3), (4) and (5) of Definition~\ref{def:genDiagDim},
    respectively.
    Moreover, it is straightforward to show that Condition~(6) holds by using the definition of ${\Phi_r}$. \qedhere
\end{proof}


\medskip
\subsection{Lower bound}

\begin{thm}  \label{thm:lowerBoundFPAlg}
    Let $X$ be a uniformly locally finite metric space.
    Then
    \begin{align*}
         \ddimn{\CfpCartan{X}}{\Cfp{\Hcal_X}}{\ell^\infty(\N, \bop{\Hcal})} \ge \asdim{X}.
    \end{align*}
\end{thm}

Our proof of Theorem~\ref{thm:lowerBoundFPAlg} uses some ideas from the proof of the
lower bound in \cite[Theorem~7.7]{LiLiaoWinter2023diagDim}, that is,
\begin{align*}
         \ddim{\ell^\infty\prn{X}}{\Cu{X}} \ge \asdim{X}.
\end{align*}
However, it differs from their argument at two points, as explained below.
Our proof does not involve groupoids and, consequently, neither the dynamic asymptotic dimension.
Instead, inspired by the proof of Theorem~6.4~in~\cite{GuentnerWillettYu2017dynAsDim}, 
we bound the diagonal dimension
directly by the asymptotic dimension. 
Moreover, Condition~(6) of the definition of the generalised diagonal dimension is essential in our proof, since it is used in order to obtain 
a uniform bound for a covering of $X$.

Before presenting the proof of the theorem, 
we prove some key lemmas in the following subsections.

\subsubsection{Setup and notation}
Suppose that $(X, \dist_X)$ is a uniformly locally finite metric space. 
For brevity, we denote the \cstarPair{} by 
    \[\big( D_A \subseteq A \big) \coloneq \big(\CfpCartan{X} \subseteq \Cfp{\Hcal_X} \big)\]
and 
\[B \coloneq \ell^\infty(\N, \bop{\Hcal}).\]
Assume that 
\[\ddimn{D_A}{A}{B}  = d < \infty.\]

Fix ${r>0}$ and set 
\begin{equation*}
    E_0 \coloneq \big\{ (x,y)\in X \times X ~ \colon ~ \distX{x}{y} \le r \big\}.
\end{equation*}
Since $X$ is uniformly locally finite, there exist disjoint sets ${S_1, S_2, \hdots, S_M \subseteq X \times X}$ such that
\[{E_0 = S_1 \sqcup S_2 \sqcup \hdots \sqcup S_M},\]
  and for each ${m\in \{1,\hdots, M \}}$ the (range and source) maps 
    \[\begin{array}{ccc}
        S_m & \rAr & X\\
        (x,y) & \longmapsto &  x
    \end{array}\] 
  and
  \[\begin{array}{ccc}
        S_m & \rAr & X\\
         (x,y) & \longmapsto &  y
    \end{array}\] 
are injective.

For each ${m \in \{1,\hdots , M\}}$ define the operator ${a_m \in \Cfp{\Hcal_X}}$, with the matrix representation
${a_m = \big( a_m(x,y) \big)_{x,y \in X}}$
given by
\begin{equation*}
        a_m (x,y) \coloneq \begin{cases}
            \id{\Hcal} , & (x,y) \in S_m \\
            0, & (x,y) \notin S_m
        \end{cases}.
\end{equation*}
    
Define the finite set
\begin{equation*}
        \Fcal \coloneq \big\{ 1_A, \, a_1 ,\hdots, \, a_M \big\} \subseteq A.
\end{equation*}
and
fix the constants 
\[\begin{array}{ccc}
        \delta \coloneq \frac{1}{2^7 (d+1)^2}\;, &
        \eta \coloneq \frac{1}{2^3 (d+1)} \;, &  \e \coloneq \frac{\delta^3}{4}
\end{array}.\]
Choose a completely positive approximation ${(F,~D_F,~\psi,~\phi)}$
witnessing 
\[\ddimn{D_A}{A}{B}=d\]
for 
${\Fcal \cup \Fcal^2 = \{ a, a^2 \colon a \in \Fcal\}}$ 
within ${{\e^2}/{81}>0}$,
where
\[F=F^{(0)} \oplus F^{(1)} \oplus \hdots \oplus F^{(d)}\]
is a finite\myhyp{}dimensional \cstar{} with the canonical diagonal 
\[D_F=D^{(0)} \oplus D^{(1)} \oplus \hdots \oplus D^{(d)}.\]
Note that, using Lemma~\ref{lem:cpPhioPsiContractive}, we can further assume that ${\phi \circ \psi}$ is contractive.

The \cstar{} $F$ is finite\myhyp{}dimensional, hence, $F$ is 
the sum of (finite\myhyp{}dimensional) matrix algebras,
\begin{equation*}
         F= \bigoplus_{k=1}^K M_{n_k}\prn{\C}.
\end{equation*}
From Condition (4) of the definition of diagonal dimension, we have 
\[ \psi(1_A) \in D_F \otimes B \subseteq F \otimes B.\]
Then we write 
\[\psi(1_A) = \sum_{k=1}^K  \psi(1_A)_{k} \in \bigoplus_{k=1}^K M_{n_k}\prn{\C} \otimes B = F \otimes B , \] 
where $\psi(1_A)_{k} \in M_{n_k}\prn{\C} \otimes B$ is the $k$\myhyp{}th entry of $\psi(1_A)$ in $\bigoplus_{k=1}^K M_{n_k}\prn{\C} \otimes B$.
    
Since $\psi(1_A) \in D_F \otimes B $, we may write 
$\psi(1_A)$ as a sum of diagonal matrices 
\[\psi(1_A)_{k} \in D_{n_k}\prn{\C} \otimes B.\] 
   
By applying the Borel functional calculus\footnote{Note that here we are using that $B$ is closed with respect to the strong operator topology.} on $B$
with $\chi_{(\delta,1]} $ 
the characteristic function on the interval $(\delta,1]$, 
we define the projection $\widehat{q} \in D_F \otimes B$ given by
\begin{equation*}
        \widehat{q}  \coloneq \chi_{(\delta,1]} (\psi(1_A) ) = \sum_k \Big(\chi_{(\delta,1]} \big( \psi(1_A)_{k} \big)\Big)
        \in  D_F \otimes B.
\end{equation*}
Let 
\[q \in  D_F \otimes B =\bigoplus_{i=0}^d D^{\prn{i}} \otimes B \] 
be the projection with
matrix entries equal to $1_B$ when the corresponding matrix entry of $\widehat{q}$ is non-zero, and $0$ elsewhere. 
We write 
\[q = (q^{\prn{0}}, q^{\prn{1}}, \hdots, q^{\prn{d}})\in  \bigoplus_{i=0}^d D^{\prn{i}} \otimes B,\] 
where ${q^{\prn{i}}\in {D^{\prn{i}} \otimes B}}$ is the $i$\myhyp{}th entry of 
${q}$ in ${\bigoplus_{i=0}^d D^{\prn{i}} \otimes B}$.

\smallskip
    
We present an example of the construction above.
\begin{example}
         Suppose that $F = M_r( \mathbb{C})$ and $D_F=D_r(\mathbb{C})$, then 
    \[ \widehat{q} = \begin{pmatrix}
        \widehat{q}_1 & 0 & 0 & \hdots &0 \\
        0 & \widehat{q}_2 & 0 & \hdots &0 \\
        0 & 0 & \widehat{q}_3 & \hdots &0 \\
        \hdots &  &  &  &\hdots \\
        0 & 0 & 0 & \hdots &\widehat{q}_r 
    \end{pmatrix} \in D_F \otimes B  = D_r(\mathbb{C}) \otimes B,\]
        and note that for each ${l\in \{1,\hdots,r\}}$ 
        we have ${\norm{\widehat{q_l}} = 0}$ or $1$,  since ${\widehat{q_l}}$ is a projection.
    Then
    \[ q =  \begin{pmatrix}
        \norm{\widehat{q}_1} 1_B & 0 & 0 & \hdots &0 \\
        0 & \norm{\widehat{q}_2} 1_B & 0 & \hdots &0 \\
        0 & 0 & \norm{\widehat{q}_3}1_B & \hdots &0 \\
        \hdots &  &  &  &\hdots \\
        0 & 0 & 0 & \hdots & \norm{\widehat{q}_r } 1_B
    \end{pmatrix} \in D_r(\mathbb{C}) \otimes B.\]
\end{example}

\smallskip
   
Note that for each $i$, the \cstar{} $F^{\prn{i}}$ is finite dimensional,
hence,
\begin{equation*}
        F^{\prn{i}} = \bigoplus_{j=1}^{N^\prn{i}} M_{n^{(i),j}}(\C),
\end{equation*}
for some $N^{\prn{i}}, n^{(i),j}\in \N$. 
We have the following inclusions ({of corner subalgebras}):
\[\begin{array}{ccc}
    q^{\prn{i}} \big( F^{(i)} \otimes B \big) q^{\prn{i}} \subseteq F^{(i)} \otimes B 
    & \text{ and } &
    q^{\prn{i}} \big( D^{(i)} \otimes B \big) q^{\prn{i}} \subseteq D^{(i)} \otimes B. 
\end{array}\] 
We identify the \cstarPair{}s 
\[\big(  q^{\prn{i}} ( D^{(i)} \otimes B ) q^{\prn{i}}  \subseteq  q^{\prn{i}} ( F^{(i)} \otimes B ) q^{\prn{i}} \big)\]
with \cstarPair{}s
\[ \Bigg( \bigoplus_{j=1}^{r^{(i)}} D_{s^{(i),j}}(B) \; \; {\subseteq}  \; \;  \bigoplus_{j=1}^{r^{(i)}} M_{s^{(i),j}}( B)   \Bigg),\]
where $r^{\prn{i}}, s^{(i),j} \in \N$ with $r^{\prn{i}} \le N^\prn{i}$ and $s^{(i),j} \le n^{(i),j}$.
    
We fix the \emph{generalised matrix units} given by
\begin{equation} \label{eq:genMatrixUnits}
       \Big( e_{k,l}^{(i),j} \Big)_{k,l=1}^{s^{(i),j}} : = \Big( e_{k,l}^{(i),j,\mathbb{C}} \otimes 1_B \Big)_{k,l=1}^{s^{(i),j}},
\end{equation}
where $e_{k,l}^{(i),j,\mathbb{C}}$ are the canonical matrix units in $ M_{s^{(i),j}}(\mathbb{C})$. 
    
For each $i,j$ we define the order zero map
\begin{equation} \label{eq:orderZeroMaps(i,j)}
        \begin{array}{ccccc} 
            \phi^{(i),j} \coloneq \phi\vert_{M_{s^{(i),j}}(B)} & \colon & M_{s^{(i),j}}(B) & \rAr & A \\
            & & b &\longmapsto & \phi(b)  
        \end{array}
\end{equation}
to be the restriction of $\phi^\prn{i}$ to $ M_{s^{(i),j}}(B) \subseteq F^{(i)} \otimes B$. 

As in \cite[Section~5]{LiLiaoWinter2023diagDim}, 
we fix the piecewise linear continuous functions 
\[{f_\delta, g_\delta \; : \; [0,1]~\rAr~\mathbb{R}}\] given by
\begin{equation}
    f_{\delta}(t) = \begin{cases}
        0 , & 0\le t\le \delta\\
        \text{linear} , & \delta < t \le 2 \delta \\
        t , & 2\delta <t \le 1
    \end{cases}
\end{equation}
and
\begin{equation}
    g_{\delta}(t) = \begin{cases}
        0 , & 0\le t\le \delta/2 \\
        \text{linear} , & \delta/2 < t \le  \delta \\
        1 , & \delta <t \le 1
    \end{cases}.
\end{equation}
    
We will use the following notation:
for each ${i\in \{0,\hdots,d\}}$, 
${j\in \{ 1,\hdots,r^{(i)}\}}$, 
and ${k,l \in \{ 1,\hdots,s^{(i),j}\}}$
set
\begin{align}
        \unMatImg{i}{j}{k,l} \coloneq \unitsimgO{f_\delta}{\phi^{(i),j}}{e_{k,l}^{(i),j}},
\end{align}
and
\begin{align}
         \gUnMatImg{i}{j}{k,l} \coloneq \unitsimgO{g_\delta}{\phi^{(i),j}}{e_{k,l}^{(i),j}},
\end{align}
where we employ the order zero functional calculus given by Corollary~\ref{cor:orderZeroCalc}.

Fix ${i\in \{0,\hdots,d\}}$ and ${j\in \{ 1,\hdots,r^{(i)}\}}$. 
With the above notation it is straightforward to prove the following result.
\begin{propo} \label{prop:unMatImgProperties} 
        For each ${k,l,m \in \{1, \hdots , r\}}$  we have:
     \begin{enumerate}
            \item[\textup{(i)}] 
            $\unMatImg{i}{j}{k,k}\in D_A$ and $\unMatImg{i}{j}{k,k} \ge 0$, 
            \item[\textup{(ii)}] 
             $\gUnMatImg{i}{j}{k,k}\in D_A$ and $\gUnMatImg{i}{j}{k,k} \ge 0$,
            \item[\textup{(iii)}] $\big(\unMatImg{i}{j}{k,l}\big)^{\ast} =\unMatImg{i}{j}{l,k}$,
            \item[\textup{(iv)}] $\big(\gUnMatImg{i}{j}{k,l}\big)^{\ast} = \gUnMatImg{i}{j}{l,k}$,
            \item[\textup{(v)}] $\unMatImg{i}{j}{k,l}
            \gUnMatImg{i}{j}{l,m } = 
            \unMatImg{i}{j}{k,m} 
            =\gUnMatImg{i}{j}{k,l}
            \unMatImg{i}{j}{l,m }$.
        \end{enumerate}
\end{propo}
\noindent
We will use these properties extensively in the following proofs without further explanation.

\smallskip
\subsubsection{Construction of partial bijections}
One can derive information from the order zero maps ${\phi^{(i),j}}$ on the structure of their images.
This result is a generalisation of \cite[Lemma~5.6]{LiLiaoWinter2023diagDim}.


\begin{lem}  \label{lem:orderZeroDecomp} 
     Fix ${i\in \{0,\hdots,d\}}$ and ${j\in \{ 1,\hdots,r^{(i)}\}}$. 
     Define the completely positive map
     \[
     \begin{array}{ccccl}
       \sigma_{k,l}^{(i),j} & \colon & A & \rAr  & \phantom{\gUnMatImg{i}{j}{l,k}} ~ A   \\
       & & a & \longmapsto & \gUnMatImg{i}{j}{l,k} ~ a  ~ \big(\gUnMatImg{i}{j}{l,k}\big)^{\ast}  
     \end{array}.
     \]
    Then $\sigma_{k,l}^{(i),j}$ restricts to a ${\ast}$-isomorphism
    \[\overline{\unMatImg{i}{j}{k,k} A \unMatImg{i}{j}{k,k}} \overset{\cong}{\longrightarrow} \overline{\unMatImg{i}{j}{l,l} A \unMatImg{i}{j}{l,l}}, \]
    and, moreover, it restricts to a ${\ast}$-isomorphism
    \[ \overline{\unMatImg{i}{j}{k,k} D_A \unMatImg{i}{j}{k,k}} \overset{\cong}{\longrightarrow} \overline{\unMatImg{i}{j}{l,l} D_A \unMatImg{i}{j}{l,l}}. \]
\end{lem}

\begin{proof}
    Theorem~\ref{thm:orderZero} implies that 
    there exist a supporting ${\ast}$\myhyp{}homomorphism $\pi$ such that 
    \[\phi^{(i),j}(a)= \phi^{(i),j}(1_{M_{s^{(i),j}}(B)}) \pi(a)\] 
    for each $a \in A$. 
    Set ${h \coloneq \phi^{(i),j}(1_{M_{s^{(i),j}}(B)}) \in (\img{\phi^{(i),j}})'}$.
    
    From the properties given by Proposition~\ref{prop:unMatImgProperties} we deduce the following:
    \begin{itemize}
        \item ${\unMatImg{i}{j}{k,k} A \unMatImg{i}{j}{k,k}}$ is closed under adjoints,
        \item $\sigma_{k,l}^{(i),j}$ is multiplicative on the $\ast$\myhyp{}algebra ${\unMatImg{i}{j}{k,k} A \unMatImg{i}{j}{k,k}}$,
        
     \item ${\sigma_{k,l}^{(i),j}\big( \unMatImg{i}{j}{k,k} A \unMatImg{i}{j}{k,k} \big) \subseteq \unMatImg{i}{j}{l,l} A \unMatImg{i}{j}{l,l}}$,
    
    \item ${\sigma_{l,k}^{(i),j} \circ \sigma_{k,l}^{(i),j}}$ is the identity map
    when restricted to ${\unMatImg{i}{j}{k,k} A \unMatImg{i}{j}{k,k}}$.
    \end{itemize}
    Therefore, by continuity, $\sigma_{k,l}^{(i),j}$ restricts to a ${\ast}$-isomorphism
    \[ \sigma_{k,l}^{(i),j}~\colon~
    \overline{\unMatImg{i}{j}{k,k} A \unMatImg{i}{j}{k,k}}  \overset{\cong}{\longrightarrow}\overline{\unMatImg{i}{j}{l,l} A \unMatImg{i}{j}{l,l}}, \]
    with inverse $\sigma_{l,k}^{(i),j}$.

    To prove the second part of the lemma, we define the continuous function ${\widetilde{g} : [0,1]~\rAr~\mathbb{R}}$, given by ${g_\delta(t) = t \widetilde{g}(t)}$ for ${t \in (0,1]}$ and ${\widetilde{g}(0)=0}$. 
    Then, using the order zero functional calculus (Corollary~\ref{cor:orderZeroCalc}), it holds
    \begin{align*}
        \gUnMatImg{i}{j}{l,k} = \phi^{(i),j}\big(e_{l,k}^{(i),j}\big) \widetilde{g}\big({\phi^{(i),j}}\big) \big(e_{k,k}^{(i),j}\big).
    \end{align*}
    Since
    ${\phi^{(i),j}(e_{k,k}^{(i),j}) = h \pi(e_{k,k}^{(i),j}) \in D_A}$, 
    we have
    \begin{align*}
       \widetilde{g}\big({\phi^{(i),j}}\big) \big(e_{k,k}^{(i),j}\big)= \widetilde{g}(h)\pi(e_{k,k}^{(i),j}) \in D_A,
    \end{align*}
    by approximating $\widetilde{g}$ with polynomials on $[0,1]$ vanishing at $0$. 
    
     From Condition~(5) of Definition~\ref{def:genDiagDim} we have
     ${\phi^{(i),j}(e_{l,k}^{(i),j})\in \normalisers{D_A}{A}}$ and then
     \begin{align*}
        \gUnMatImg{i}{j}{l,k} = \phi^{(i),j}\big(e_{l,k}^{(i),j}\big) \widetilde{g}\big({\phi^{(i),j}}\big) \big(e_{k,k}^{(i),j}\big) \in \normalisers{D_A}{A}.
     \end{align*}
     Therefore,  $\sigma_{k,l}^{(i),j}$ maps $D_A$ into $D_A$. 
     Moreover, for $a \in D_A$, we obtain
     \begin{align*}
         \sigma_{k,l}^{(i),j}\big(  \unMatImg{i}{j}{k,k} a \unMatImg{i}{j}{k,k} \big) 
        & = \gUnMatImg{i}{j}{l,k} \unMatImg{i}{j}{k,k} a \unMatImg{i}{j}{k,k}  \gUnMatImg{i}{j}{k,l} =  \unMatImg{i}{j}{l,l} \gUnMatImg{i}{j}{l,k} a \gUnMatImg{i}{j}{k,l}   \unMatImg{i}{j}{l,l}.
     \end{align*}
     Hence,
     \[\sigma_{k,l}^{(i),j}\big( \unMatImg{i}{j}{k,k} a \unMatImg{i}{j}{k,k} \big) \in \unMatImg{i}{j}{l,l} D_A \unMatImg{i}{j}{l,l}.\]
    This proves that $\sigma_{k,l}^{(i),j}$ restricts to a ${\ast}$-isomorphism 
     \[ \overline{\unMatImg{i}{j}{k,k} D_A \unMatImg{i}{j}{k,k}} \overset{\cong}{\longrightarrow} \overline{\unMatImg{i}{j}{l,l} D_A \unMatImg{i}{j}{l,l}}. \qedhere \]
\end{proof}

 Recall that 
 ${\eta = \frac{1}{2^3 (d+1)}}$
 and define the sets
\begin{equation*}
        {U^{(i),j}_k} \coloneq \Big\{ x \in X  \colon  \normbig{
        \unMatImg{i}{j}{k,k}
         \indF_x \unMatImg{i}{j}{k,k} }> \eta^2 \Big\},
\end{equation*}
for each ${i \in \{ 0, \hdots, d\}}$, ${j\in \{ 1,\hdots,r^{(i)}\}}$ and ${k\in \{1,\hdots,s^{(i),j}\}}$.
    
We define bijections 
\[{{\overline{\sigma}_{k,l}^{(i),j}}~\colon~ {U^{(i),j}_k}~\rAr~U^{(i),j}_l }\]
induced by ${{\sigma}_{k,l}^{(i),j}}$.

\begin{lem} \label{lem:homeoDecomp}
     Fix ${i\in \{0,\hdots,d\}}$, ${j\in \{ 1,\hdots,r^{(i)}\}}$
     and ${k,l\in \{1,\hdots,r\}}$.
        Then, for each ${x \in{U^{(i),j}_k}}$, there exists a ${y_x \in{U^{(i),j}_l}}$ such that 
     the operator 
     \begin{equation*}
         \sigma_{k,l}^{(i),j}(\unMatImg{i}{j}{k,k} \indF_x \unMatImg{i}{j}{k,k}) \in D_A 
     \end{equation*}
     is supported in $\{y_x\}$.
    Moreover, the assignment 
    \[
     \begin{array}{ccccl}
      {\overline{\sigma}_{k,l}^{(i),j}} & \colon & {U^{(i),j}_k}  & \rAr  &  U^{(i),j}_l\\
      & & x & \longmapsto & y_x 
     \end{array}
     \] 
    is a bijection, with 
    inverse 
    \[
     \begin{array}{ccccl}
      \left({\overline{\sigma}_{k,l}^{(i),j}} \right)^{-1}= {\overline{\sigma}_{l,k}^{(i),j}} & \colon & {U^{(i),j}_l}  & \rAr  &  U^{(i),j}_k
     \end{array}.
     \] 
\end{lem}

\begin{proof}[Sketch of proof]
     Let ${x \in U^{(i),j}_k}$. Set ${\xi \coloneq \sigma_{k,l}^{(i),j}(\unMatImg{i}{j}{k,k} \indF_x \unMatImg{i}{j}{k,k}) \in D_A }$. 
    Since 
    \begin{align*}
        \normbig{\xi} = \normbig{ \sigma_{k,l}^{(i),j}(\unMatImg{i}{j}{k,k} \indF_x \unMatImg{i}{j}{k,k}) } =  \normbig{ \unMatImg{i}{j}{k,k} \indF_x \unMatImg{i}{j}{k,k}} >\eta^2 ,
    \end{align*}
    there exists ${y \in X}$ such that 
    for ${\xi_y \coloneq \xi(y) \in \bop{\Hcal}}$ it holds
     $ {\norm{\xi_y} > \eta^2}$.

     Using the properties of Proposition~\ref{prop:unMatImgProperties} 
     and the results of Lemma~\ref{lem:orderZeroDecomp}, more precisely, 
     that $\sigma_{k,l}^{(i),j}$ is a ${\ast}$-isomorphism
    \[{\overline{\unMatImg{i}{j}{k,k} A \unMatImg{i}{j}{k,k}} \overset{\cong}{\longrightarrow} \overline{\unMatImg{i}{j}{l,l} A \unMatImg{i}{j}{l,l}} }\]
    and restricts to a ${\ast}$-isomorphism
    \[\overline{\unMatImg{i}{j}{k,k} D_A \unMatImg{i}{j}{k,k}} \overset{\cong}{\longrightarrow} \overline{\unMatImg{i}{j}{l,l} D_A \unMatImg{i}{j}{l,l}},\]
    it is straightforward to prove that there exists a unique ${y \in X}$ such that
    \begin{align*}
            \xi(z) =\begin{cases}
                \xi_y , & z=y\\
                0, &  \text{else}
            \end{cases},
        \end{align*}
    for each $z \in X$,
    and 
    \[{y \in{U^{(i),j}_l}}.\]
    For each ${k,l\in\{1,\hdots,r\}}$ we define the functions
    \begin{equation}
       {\overline{\sigma}_{k,l}^{(i),j}} \colon{U^{(i),j}_k}~\rAr~U^{(i),j}_l,
    \end{equation}
    given by 
    \[{{\overline{\sigma}_{k,l}^{(i),j}}(x) = y},\]
    where ${y \in X}$ is the unique $y$ defined as above. 
    Using Proposition~\ref{prop:unMatImgProperties} 
    and Lemma~\ref{lem:orderZeroDecomp} we can show that
    \begin{equation*}
        {\overline{\sigma}_{l,k}^{(i),j}} \circ {\overline{\sigma}_{k,l}^{(i),j}} = \id{U^{(i),j}_k} \colon{U^{(i),j}_k}~\rAr~U^{(i),j}_k.
    \end{equation*}
    
    Combining the above we obtain that
    \begin{equation*}
        {\overline{\sigma}_{k,l}^{(i),j}} \colon{U^{(i),j}_k} \xlongrightarrow{\simeq}{U^{(i),j}_l}
    \end{equation*}
    is a bijection with inverse given by ${\overline{\sigma}_{l,k}^{(i),j}} \colon{U^{(i),j}_l}~\rAr~U^{(i),j}_k$. \qedhere
\end{proof}

\smallskip
\subsubsection{Construction of the covering}
    We prove that the collection
    \begin{equation*}
        {\mathcal{C}} \coloneq \Big\{ {U}^{(i),j}_k \colon i = 0,\hdots,d, \; j= 1,\hdots,r^{(i)},\; k=1,\hdots,s^{(i),j} \Big\}
    \end{equation*}
    is a cover of $X$.

    \begin{lem} \label{lem:coverAsdim}
        ${\mathcal{C}}$ forms a cover of $X$.
    \end{lem}
    \begin{proof}
        Let $x \in X$. 
        Using
        \begin{align*}
            \normBig{   \sum_{j =1 }^{r^{(i)}} \sum_{k =1 }^{s^{(i),j}} \unMatImg{i}{j}{k,k} - 
            {\phi^{(i),j}}\big({e_{k,k}^{(i),j}} \big)} 
            & \le \normBig{f_\delta - [z\mapsto z]}_{\cont{[0,1]}} \normBig{  \sum_{j,k } {\phi^{(i),j}}\big({e_{k,k}^{(i),j}} \big)}
            \le \delta,
        \end{align*}  
        we obtain
        \begin{align*}
            \sum_{i=0}^d \normBig{ \sum_{j =1 }^{r^{(i)}} \sum_{k =1 }^{s^{(i),j}}
             \unMatImg{i}{j}{k,k}  \indF_x }
             & \ge \sum_{i=0}^d \normBig{ \sum_{j,k }
            {\phi^{(i),j}}\big({e_{k,k}^{(i),j}} \big)\indF_x }  - (d+1)\delta.
        \end{align*}
        
    Then 
    \begin{align*}
            \sum_{i=0}^d  \normBig{ \sum_{j =1 }^{r^{(i)}} \sum_{k =1 }^{s^{(i),j}}
           \unMatImg{i}{j}{k,k} \indF_x }
             & \ge  \normBig{ \sum_{i, j,k }
            {\phi^{(i),j}}\big({e_{k,k}^{(i),j}} \big)\indF_x }  - (d+1)\delta\\
            & =  \normBig{ \sum_{i }
            {\phi^{(i)}}\big( q^{\prn{i}} 1_{F^{(i)}\otimes B} q^{\prn{i}}\big)\indF_x }  - (d+1)\delta\\
             & =  \normBig{ \phi \big( q^2 \big)\indF_x }  - (d+1)\delta\\
              & \ge  \normBig{ \phi \big( q \psi(1_A) q\big)\indF_x }  - (d+1)\delta
        \end{align*}
        where for the last inequality we have used that $\phi$ is positive and $\psi$ is contractive. 
        
       Using the definition of $q$ and $\widehat{q}$, we observe that
        \begin{align*}
            \norm{\phi \circ \psi (1_A) - \phi (q \psi (1_A) q) } 
           & \le \norm{\phi}\norm{ \psi (1_A) - q  \psi (1_A) q }  \\
           & \le (d+1) \norm{ \psi (1_A) - q  \psi (1_A)} \\
            & \le (d+1) ( \delta
            + \norm{ q \; \widehat{q}\;  \psi (1_A) - 
            q \psi (1_A)})  \\
            & \le 2(d+1) \delta.
        \end{align*}
         
        Combining the above, we obtain
        \begin{align*}
           \sum_{i=0}^d \Big\lVert \sum_{j =1 }^{r^{(i)}} \sum_{k =1 }^{s^{(i),j}}
             \unMatImg{i}{j}{k,k} \indF_x \Big\rVert
             & \ge  \normBig{ \phi \left( \psi(1_A) \right)\indF_x } - 3(d+1) \delta .
        \end{align*}
         Recall that 
         ${\delta = \frac{1}{2^7 (d+1)^2}}$ and  ${\e = \frac{\delta^3}{4}}$. 
         Since ${1_A\in \Fcal}$ is approximated by ${\phi \circ \psi (1_A)}$, we have
         \begin{align*}
           \sum_{i=0}^d \Big\lVert \sum_{j =1 }^{r^{(i)}} \sum_{k =1 }^{s^{(i),j}}
             \unMatImg{i}{j}{k,k} \indF_x \Big\rVert
             & \ge  \normBig{ 1_A\indF_x } - \e - 3(d+1) \delta  
              > \frac{3}{4}.
        \end{align*}
              
        Then there exists ${i_x \in \{0,\hdots,d\}}$ such that
        \begin{align*}
         \norm{ \sum_{j =1 }^{r^{(i_x)}} \sum_{k =1 }^{s^{(i_x),j}}
             \unMatImg{i_x}{j}{k,k} \indF_x }
             > \frac{3}{4(d+1)} > \eta.
        \end{align*}
        Since $\phi^{(i_x)}$ is an order zero map, we have that the summands above are pairwise orthogonal 
        and hence there exists ${j_x \in \{1,\hdots,r^{(i_x)}\}}$ and ${k_x \in \{1,\hdots,s^{(i_x),j_x}\}}$ such that
     \[{\norm{ \unMatImg{i_x}{j_x}{k_x,k_x} \, \indF_x }
             > \eta}.\]
    Therefore, ${x \in {U}_{k_x}^{(i_x),j_x}}$. This proves that the collection ${{\mathcal{C}}}$ forms a cover for $X$. \qedhere
     \end{proof}

 For each $i\in \{0,\hdots,d\}$ set  
    \begin{equation*}
        U^{(i)} \coloneq \bigcup_{j=1}^{r^{(i)}} \bigcup_{k=1}^{s^{(i),j}}  U^{(i),j}_k . 
    \end{equation*}
    Define the equivalence relation\footnote{This was inspired by the equivalence relations defined in 
   ~\cite[Theorem~6.4]{GuentnerWillettYu2017dynAsDim} and~\cite[Proposition~6.6]{LiLiaoWinter2023diagDim}.} 
   on $U^{(i)}$
    given by:
    \begin{equation*}
        x \sim_i y
    \end{equation*}
    if and only if there are 
    $x_1,\hdots,x_\kappa \in U^{(i)}$  
    such that:
    \begin{itemize}
        \item $x_1 = x,$
        \item $ x_k = y,$
        \item ${\distX{x_l}{x_{l+1}} \le r}$ for each $l \in \{1,\hdots,\kappa-1\}$.
    \end{itemize}
  
    We claim that in order to bound the asymptotic dimension, 
    it suffices to prove that the cardinalities
    of the equivalence classes ${[x]_{i}}$ are uniformly bounded.
    
    \begin{lem}\label{lem:LowerBoundProof1}
    If
    \begin{equation*}
        S\coloneq \max_{i=0,\hdots,d} \bigg(\sup_{x \in U^{(i)}} \big{\lvert}{ [ x ]_{i}}   \big{\rvert} \bigg) < \infty ,
    \end{equation*}
    then
    \[\asdim{X} \le d.\]
    \end{lem}
    
    \begin{proof}
     For each ${i\in \{0,\hdots,d\}}$ define the set of all equivalence classes
    \begin{equation*}
        \Big\{ V^{(i),n} \colon n \in N_i  \Big\} \coloneq \Big\{ [ x ]_{i} \colon  x \in U^{\prn{i}} \Big\}  ,
    \end{equation*}
    where $N_i$ is an enumeration of the classes ${[\,\cdot\,]_{i}}$.
    Moreover, define the collections
    \begin{equation*}
        \Vcal^{(i)} \coloneq  \Big\{ V^{(i),n} \colon n \in N_i  \Big\}
    \end{equation*}
    and set
    \begin{equation*}
        \Vcal \coloneq  \Vcal^{(0)} \sqcup \hdots \sqcup  \Vcal^{(d)}.
    \end{equation*}
    We will show that $\Vcal$ is a covering that bounds the asymptotic dimension. 

    From Lemma~\ref{lem:coverAsdim} we have that ${\mathcal{C}}$ covers $X$, 
  therefore $\Vcal$ covers $X$.

    Let ${x, y \in V^{(i),n}}$ for some ${i\in \{0,\hdots,d\}}$ and ${n \in N_i}$. 
    Then ${x \sim_i y}$ and, since ${S < \infty }$, we have
    ${\distX{x}{y} \le S r}$.
    Thus, 
    \[\diam{V^{(i),n}} \le S r,\]
    for each ${i\in \{0,\hdots,d\}}$ and ${n \in N_i}$.
    In other words, ${\Vcal}$ is uniformly bounded.
      
    To show that $\Vcal$ is decomposed into $r$-separated collections, 
    we fix ${i\in \{0,\hdots,d\}}$ and ${n, n' \in N_i}$ with ${n \ne n'}$. 
    Suppose that ${x, y \in X}$, satisfying 
    \[(x,y) \in  \big(V^{(i),n}  \times  V^{(i),n'}  \big) \]
    and
    \[ {\distX{x}{y} \le r}. \]
    Note that ${\distX{x}{y} \le r}$
    implies
    ${[ x ]_{i} = [ y ]_{i}}$.
    This is a contradiction because ${n \ne n'}$. Therefore,
    \[ \Vcal^{(i)} =  \Big\{ V^{(i),n} \colon n \in N_i  \Big\} \]
    is an $r$-separated collection.

     Combining the above we obtain 
        ${\asdim{X} \le d}$.  \qedhere    
    \end{proof}
    

\smallskip
\subsubsection{Some technical definitions}
    Inspired by Remark~2.2.(ii) in \cite{LiLiaoWinter2023diagDim}, 
    we define completely positive contractive maps $\widehat{\psi}$ and $\widehat{\phi}$ 
    such that ${\widehat{\phi} \circ \widehat{\psi}}$ approximates the identity map on $\Fcal$.
    
    Set ${\widehat{F}}$ to be the \cstarSub{} of ${F \otimes B}$ generated by 
    \[\psi(1_A) \big( F \otimes B \big) \psi(1_A) .\]
    Define the functions ${\zeta, \zeta' \in \cont{[0,1]}}$ given by 
    \[
        \zeta(z) \coloneq \begin{cases}
            \frac{\e}{9 \sqrt{2 (d+1)}} , & 0 \le z \le \frac{\e^2}{81 \cdot 2 (d+1)}\\
            \sqrt{z}, & \frac{\e^2}{81 \cdot 2 (d+1)} < z \le 1
        \end{cases} \]
      and
      \[\begin{array}{cc}
           \zeta'(z) \coloneq \frac{1}{\zeta(z)}, &  0 \le z  \le 1 
      \end{array}.\]
    Using the continuous functional calculus, define ${p, p' \in \widehat{F}}$ by 
    \[\begin{array}{ccc}
        p \coloneq \zeta(\psi(1_A)) 
        & \text{ and }  &
     p' \coloneq \zeta'(\psi(1_A)),
    \end{array}\]
     and
     observe that $p$ and $p'$ are positive.
   
     We define the following maps:
    \[\begin{array}{ccccc}
        \widehat{\psi}& \colon & A
        &~\rAr~& \widehat{F}\\
       & & x &~\longmapsto~& p' \psi(x) p' \end{array},\]
       and
       \[\begin{array}{ccccc}
        \widehat{\phi}& \colon & \widehat{F} 
        &\longrightarrow & A \\
        & & x &~\longmapsto~& \frac{\phi( p  x  p )}{1+\e^2/81}   
    \end{array}.\]
    
    Using the above definitions we obtain the following result.

    \begin{propo}  \label{prop:capCpcMaps}
    The  maps 
    \[ A \overset{\widehat{\psi}}{\longrightarrow} \widehat{F} \overset{\widehat{\phi}}{\longrightarrow} A,\]
    are
    completely positive contractive maps
    such that:
        \begin{enumerate} \renewcommand{\labelenumi}{(\roman{enumi})}
            \item for each $a \in A$, it holds
            \begin{equation*}
             \phi \circ \psi (a) = \frac{1}{1+\e^2/81} \widehat{\phi} \circ \widehat{\psi} (a),
            \end{equation*}
            \item 
                 for each $a \in \Fcal \cup \Fcal^2$, it holds
            \begin{equation*}
        \normBig{\widehat{\phi} \circ \widehat{\psi}(a) - a} < \e^2/27,
    \end{equation*}
            \item for each $a \in \Fcal$, $b \in \widehat{F}$, with $\norm{b} \le 1$, it holds
            \begin{equation*} \label{ineq:capCpcMapsBound}
        \normBig{ \widehat{\phi} \big( \widehat{\psi}(a) b \big) - \widehat{\phi}\big(\widehat{\psi}(a)\big) \widehat{\phi}\big(b\big)} 
        < 6 \Big(\frac{\e^2}{81}\Big)^{1/2}  < \e.
    \end{equation*}
        \end{enumerate}
    \end{propo}

    It should be stressed that in this case $\widehat{\phi}$ is no longer a sum of order zero maps. 

   It is straightforward to show Proposition~\ref{prop:capCpcMaps}.(i) and (ii) by 
   using the definitions of ${\widehat{\psi}}$ and ${\widehat{\phi}}$, and the approximation property of ${\phi \circ \psi}$.
   Proposition~\ref{prop:capCpcMaps}.(iii) follows from Proposition~\ref{prop:capCpcMaps}.(ii) 
   and \cite[Lemma~3.5]{KirchbergWinter2004CoveringDimQuasidiagonality}. 
   Detailed computations can be found in \cite[Lemma~4.9]{Kitsios2025PhDThesis}.


    \smallskip
    \subsubsection{Conclusion of the proof}
    We now have all the ingredients to prove Theorem~\ref{thm:lowerBoundFPAlg}.

    \begin{proof}[Proof of Theorem~\ref{thm:lowerBoundFPAlg}]
    By Lemma~\ref{lem:LowerBoundProof1}, in order to prove the theorem, it suffices to show that 
     \begin{equation*}
        S = \max_{i=0,\hdots,d} \bigg(\sup_{x \in U^{(i)}} \big{\lvert}{ [ w ]_{i}} \big{\rvert} \bigg) < \infty .
    \end{equation*}

        Fix ${i \in \{0,\hdots, d\}}$.
      Let ${w \in U^{\prn{i}}}$ and without loss of generality assume that ${w \in U^{(i),j_0}_{k_0}}$,
      for some ${j_0}$ and ${k_0}$. 
        Suppose that ${\widetilde{w} \in U^{\prn{i}}}$ is such that ${w \sim_i \widetilde{w}}$. 
        Then there exist 
        \[{w_1, \hdots, w_\kappa \in U^{\prn{i}}}\]
        with 
    \begin{itemize}
        \item $w_1 = w,$
        \item $w_\kappa = \widetilde{w},$
        \item ${\distX{w_l}{w_{l+1}} \le r}$ for each ${l \in \{1,\hdots,\kappa-1\}}$.
    \end{itemize}

    Fix ${l \in \{1,\hdots,\kappa-1\}}$ and set 
    ${z=w_l}$ and ${z'=w_{l+1}}$. 
    Without loss of generality, assume that
    ${z \in U^{(i),j}_k}$ and ${z' \in U^{(i),j'}_{k'}}$ for some  ${j, j'}$ and ${k, k'}$. 
        Set 
        \begin{equation*}
            x \coloneq {{\overline{\sigma}}_{k,1}^{(i),j}}(z)  \in U^{(i),j}_1 
        \end{equation*}
        and
        \begin{equation*}
        y \coloneq {\overline{\sigma}_{k',1}^{(i),j'}}(z')  \in U^{(i),j'}_1.
        \end{equation*}
       Note that
        \[{{\overline{\sigma}}_{1,k}^{(i),j}} \big(x \big) = {{\overline{\sigma}}_{1,k}^{(i),j}} \big({{\overline{\sigma}}_{k,1}^{(i),j}}(z) \big) = z\]
        and
        \[{\overline{\sigma}_{1,k'}^{(i),j'}}\big(y \big) = {\overline{\sigma}_{1,k'}^{(i),j'}}\big({\overline{\sigma}_{k',1}^{(i),j'}}(z') \big)
        = z'. \]
        
        We will show that 
        $j = j'$ and $x=y$, and, in particular, it holds
        \begin{equation*}
            z' = {\overline{\sigma}_{1,k'}^{(i),j}}\big({{\overline{\sigma}}_{k,1}^{(i),j}}(z) \big).
        \end{equation*}

        We define the operators
        \begin{equation*}
            T_x \coloneq \indF_x \, \gUnMatImg{i}{j}{1,k} \, \indF_{z} \, \gUnMatImg{i}{j}{k,1} \, \indF_x \in D_A
        \end{equation*}
        and 
        \begin{equation*}
        \widetilde{T}_{z} \coloneq \indF_{z} \, \gUnMatImg{i}{j}{k,1} T_x \gUnMatImg{i}{j}{1,k} \, \indF_{z} \in D_A,
        \end{equation*}
        which are supported on $\{x\}$ and $\{z\}$, respectively.
        Since ${z \in U^{(i),j}_k}$, it holds
         \begin{equation*}
        \normbig{\unMatImg{i}{j}{k,k}  \widetilde{T}_{z} \unMatImg{i}{j}{k,k} }  > \eta^2.
    \end{equation*}
         Moreover, using the above estimate and the definition of ${\widehat{\phi}}$ we see that
        \begin{equation*}
            \norm{ \widehat{\phi}\Big(e_{k,k}^{(i),j} \Big)  \widetilde{T}_{z} \unMatImg{i}{j}{k,k}}  >  \delta^2/2.
        \end{equation*}
        Set 
       \begin{equation*}
           \widehat{T} \coloneq  \widehat{\phi}\left(e_{k,k}^{(i),j} \right)   \widetilde{T}_{z} \unMatImg{i}{j}{k,k} \in D_A,
       \end{equation*}
       and observe that $\widehat{T}$ is supported on $\{z\}$, 
       since $\phi$ maps diagonal elements into $D_A$.
        
        Since ${\distX{z'}{z} \le r}$, there exists a unique ${m \in \{1,\hdots,M\}}$ such that
     \[{a_m(z',z) = \id{\Hcal}}.\] 
        Then, using the matrix representation of operators in $A$, 
    it is straightforward to see that
    \begin{equation*}
                \widehat{T}(z,z) 
                = (a_m \widehat{T} a_m^\ast) (z',z') = (\indF_{z'} a_m \widehat{T} a_m^\ast) (z',z'),
    \end{equation*}
    where ${(\indF_{z'} a_m \widehat{T} a_m^\ast) (z',z')}$ is the ${(z',z')}$\myhyp{}entry in the matrix representation of the operator ${\indF_{z'} a_m \widehat{T} a_m^\ast\in A}$.

        By combining the above
        \begin{align*}
          \frac{\delta^2}{2}
          & <  \normBig{\widehat{T}} 
          = \normBig{\widehat{T}(z,z)}
          = \normBig{(\indF_{z'} a_m \widehat{T} a_m^\ast) (z',z')}
          \le \normBig{\indF_{z'} a_m \widehat{T} a_m^\ast }.
        \end{align*}

        Since ${a_m \in \Fcal}$, using Proposition~\ref{prop:capCpcMaps}.(ii), we obtain
         \begin{align*}
          \frac{\delta^2}{2}
          < \normBig{\indF_{z'} \widehat{\phi} \big(\widehat{\psi}(a_m) \big)  \widehat{T} a_m^\ast } + \frac{\e^2}{27}.
        \end{align*}
        Then, by the definition of $\widehat{T}$ and Proposition~\ref{prop:capCpcMaps}.(iii), we have
        \begin{align*}
          \frac{\delta^2}{2}
           & < \normBig{\indF_{z'} \widehat{\phi} \big(\widehat{\psi}(a_m) \big)  \widehat{T} a_m^\ast } + \frac{\e^2}{27}\\
           & = \normBig{\indF_{z'}  \, \widehat{\phi} \big(\widehat{\psi}(a_m) \big) \widehat{\phi}\big(e_{k,k}^{(i),j} \big)  \widetilde{T}_{z} \unMatImg{i}{j}{k,k} a_m^\ast} + \frac{\e^2}{27}\\
           & \le \normBig{\indF_{z'} \, \widehat{\phi}\big(\widehat{\psi}(a_m) e_{k,k}^{(i),j} \big)   \widetilde{T}_{z} \unMatImg{i}{j}{k,k} a_m^\ast} + \frac{\e^2}{27} +\e,
        \end{align*}
        and, using the definition of
        $\widehat{\psi}$ and $\widehat{\phi}$,
        we obtain
        \begin{align*}
          \frac{\delta^2}{2}
           & \le   \normBig{\indF_{z'} \, {\phi}\big({\psi}(a_m) e_{k,k}^{(i),j} \big)   \widetilde{T}_{z} \unMatImg{i}{j}{k,k} a_m^\ast} + \frac{\e^2}{27} +\e.
        \end{align*}
  
        Define 
        \[{h^{(i)} \coloneq \phi^{(i)} \big( 1_{F^{(i)} \otimes B} \big)}\]
        and suppose that
         \[\pi^{(i)} \colon F^{(i)} \otimes B~\rAr~ \bop{\ell^2(X,\Hcal)}\]
         is a supporting ${\ast}$\myhyp{}homomorphism of the 
         order zero map $\phi^{(i)}$, that is,
        \[{\phi^{(i)}(b) = h^{(i)} \pi^{(i)}(b)},\]
        for each ${b \in F^{(i)} \otimes B}$.
    Condition~(6) in the definition of generalised diagonal dimension implies that
    \begin{equation*}
        \indF_{z' }= \pi^{(i)}\big(e^{{(i)},{j'}}_{k',k'} \big) \; \indF_{z' } \;\pi^{(i)}\big(e^{{(i)},{j'}}_{k',k'} \big) .
    \end{equation*}

        Note that
        ${\unMatImg{i}{j'}{1,k'} \indF_{z'} \unMatImg{i}{j'}{k',1}}$
        is supported in $\{y\}$,
        since ${z' = \overline{\sigma}_{1,k'}^{(i),j'}\big(y \big)}$.
        Then, using the definition of ${\unMatImg{i}{j'}{1,k'}}$
    and the definition of $\pi^{(i)}$, we obtain that the operator
    \begin{equation*}
        P_y \coloneq \pi^{(i)}\big(e^{\prn{i},j'}_{1,k'} \big) \; \indF_{z'} \;\pi^{(i)}\big(e^{\prn{i},j'}_{k',1} \big) \in D_A 
    \end{equation*}
    is supported in $\{y\}$.
    We have
    \begin{align*}
        \pi^{(i)}\big(e^{\prn{i},j'}_{k',1} \big)  P_y \pi^{(i)}\big(e^{\prn{i},j'}_{1,k'} \big)
        = \pi^{(i)}\big(e^{\prn{i},j'}_{k',k'} \big) \; \indF_{z' } \;\pi^{(i)}\big(e^{\prn{i},j'}_{k',k'} \big) 
        = 
        \indF_{z'}.
    \end{align*}

        By replacing the above in the previous inequality, we obtain
        \begin{align} \label{ineq:xNEyIneq}
          \frac{\delta^2}{2}
           <  \normBig{\pi^{(i)}\big(e^{\prn{i},j'}_{k',1} \big)  P_y \pi^{(i)}\big(e^{\prn{i},j'}_{1,k'} \big) {\phi}\big({\psi}(a_m) e_{k,k}^{(i),j} \big)   \widetilde{T}_{z} \unMatImg{i}{j}{k,k} a_m^\ast} + \frac{\e^2}{27} +\e.
        \end{align}

        Since $\phi^{(i)}$ is an order zero map with $\pi^{\prn{i}}$ a supporting $\ast$\myhyp{}homomorphism,
        ${j \ne j'}$ 
        implies that the first summand above is zero, and then
        \[\frac{\delta^2}{2}
          < \frac{\e^2}{27} +\e,\]
        which is a contradiction. 
        Therefore, $j=j'$.

        Assume now that ${x \ne y}$ and
        set 
        \[v \coloneq e^{{(i)},{j}}_{1,k'} \psi(a_m) e^{{(i)},{j}}_{k,1} \in F^{\prn{i}} \otimes B.\] 
        Observe that, by its definition,
        $v$ is of the form ${v = e^{{(i)},{j},\C}_{1,1} \otimes b}$ for some ${b \in B}$.
        
        Then, 
        using that ${\pi^{\prn{i}}}$ is a supporting $\ast$\myhyp{}homomorphism of $\phi^{(i)}$, 
        and that 
        \[\gUnMatImg{i}{j}{k,1} =g_\delta\big(h^{(i)} \big) \pi^{(i)}\big(e^{{(i)},{j}}_{k,1}\big),\]
        we obtain
        \begin{align*}
            \pi^{(i)}\big(e^{\prn{i},j}_{1,k'} \big) {\phi}\big({\psi}(a_m) e_{k,k}^{(i),j} \big)   \widetilde{T}_{z}
            & = \pi^{(i)}\big(e^{\prn{i},j}_{1,k'} \big) {\phi}\big({\psi}(a_m) e_{k,k}^{(i),j} \big)    \gUnMatImg{i}{j}{k,1} T_x \gUnMatImg{i}{j}{1,k}  \indF_{z}\\
            & =  {\phi}\big( e^{\prn{i},j}_{1,k'} {\psi}(a_m) e_{k,k}^{(i),j} e^{\prn{i},j}_{k,1} \big)   g_\delta\big(h^{(i)}\big) T_x \gUnMatImg{i}{j}{1,k}  \indF_{z}\\
           &  =  {\phi}\big( v \big)   g_\delta\big(h^{(i)}\big) T_x \gUnMatImg{i}{j}{1,k}  \indF_{z}.
        \end{align*}
    
        The following:
        \begin{itemize}
            \item ${v \in D^{\prn{i}} \otimes B}$,
            \item ${h^{(i)} = \phi^{(i)} \big( 1_{F^{(i)} \otimes B} \big)}$,
            \item $\phi$ maps the diagonal into $D_A$,
        \end{itemize}
        imply that 
         \[{{\phi}\big( v \big)   g_\delta\big(h^{(i)}\big) \in D_A}.\]
        Recall that ${P_y \in D_A}$ is supported on $\{y\}$ and ${T_x \in D_A}$ is supported on $\{x\}$. 
        Then, using that ${x \ne y}$, we obtain
        \begin{align*}
            P_y \pi^{(i)}\big(e^{\prn{i},j}_{1,k'} \big) {\phi}\big({\psi}(a_m) e_{k,k}^{(i),j} \big)   \widetilde{T}_{z}
            = P_y {\phi}\big( v \big)   g_\delta\big(h^{(i)}\big) T_x \gUnMatImg{i}{j}{1,k}  \indF_{z} = 0.
        \end{align*}
          
         Therefore, by combining the above and
         Inequality~(\ref{ineq:xNEyIneq}), we have
         \begin{align*}
          \frac{\delta^2}{2}
           & <  \normBig{\pi^{(i)}\big(e^{\prn{i},j'}_{k',1} \big)  P_y \pi^{(i)}\big(e^{\prn{i},j'}_{1,k'} \big) {\phi}\big({\psi}(a_m) e_{k,k}^{(i),j} \big)   \widetilde{T}_{z} \unMatImg{i}{j}{k,k} a_m^\ast} + \frac{\e^2}{27} +\e
           = \frac{\e^2}{27} +\e,
        \end{align*}
        which is a contradiction. 
        Thus, $x=y$.
  
         Therefore, we obtain
         \begin{equation} \label{eq:LowerBoundProof1}
                w_{l+1} = z' 
         ={\overline{\sigma}_{1,k'}^{(i),j}}\big(y \big)
        = {\overline{\sigma}_{1,k'}^{(i),j}}\big(x \big)
        = {\overline{\sigma}_{1,k'}^{(i),j}}\big({{\overline{\sigma}}_{k,1}^{(i),j}}(z) \big)
        = {\overline{\sigma}_{1,k'}^{(i),j}}\big({{\overline{\sigma}}_{k,1}^{(i),j}}(w_l) \big). 
         \end{equation}

    Recall that ${w = w_1, w_2, \hdots, \widetilde{w}=w_\kappa \in U^{\prn{i}}}$ are such that ${\distX{w_l}{w_{l+1}} \le r}$ for each ${l}$. 
    Then, by applying the previous (Equation~\eqref{eq:LowerBoundProof1}) recursively to the pairs ${(w_{m}, w_{m+1})}$ for each ${m=1,\hdots,\kappa-1}$, 
    we obtain 
    \begin{equation*}
        \widetilde{w} = w_\kappa \in \Big\{{{\overline{\sigma}}_{1,k'}^{(i),j_0}}\big({\overline{\sigma}}_{k_0,1}^{(i),j_0}(w) \big) \colon k' = 1,\hdots,s^{(i),j_0}  \Big\}.
    \end{equation*}
    This implies 
        \[ \big{\lvert}{ [ w ]_{i}} \big{\rvert}  \le s^{(i),j_0} \]
        and, therefore,
        \begin{equation*}
            S = \max_{i=0,\hdots,d} \bigg(\sup_{x \in U^{(i)}} \big{\lvert}{ [ w ]_{i}} \big{\rvert} \bigg) \le \max_{ i=0,\hdots,d  } \bigg(\sup_{j=1,\hdots,r^{(i)}} \abs{ s^{(i),j} }\bigg)  < \infty . \qedhere
        \end{equation*}
\end{proof}

\begin{rmk}
It should be noted that in the last part of the proof, that is, to show ${S<\infty}$,
we used Condition~(6) of the generalised diagonal dimension
(Definition~\ref{def:genDiagDim}).
\end{rmk}


\medskip
\subsection{Concluding remarks}

\subsubsection{Generalised diagonal dimension of Roe algebras}

We expect that one can extract the asymptotic dimension of a uniformly locally finite metric space $X$ from its Roe algebra. 
More specifically, 
by using similar techniques as in Theorems \ref{thm:upperBoundFPAlg} and \ref{thm:lowerBoundFPAlg},
one should be able to prove
\begin{equation*}
         \ddimn{\CroeCartan{X}}{\Croe{\ell^2(X, \Hcal)}}{\ell^\infty(\N, \kop{\Hcal})} = \asdim{X}.
\end{equation*}
One technical obstruction that appears
is that in our proof of Theorem \ref{thm:lowerBoundFPAlg}
we used that ${\ell^\infty(\N, \bop{\Hcal})}$ is closed in the strong operator topology,
whereas ${\ell^\infty(\N, \kop{\Hcal})}$ is not.
To overcome this, one could use that
the \cstar{} of finite\myhyp{}propagation operators on $X$
is the multiplier algebra of the Roe algebra ${\Croe{X}}$.
By overcoming this obstruction, one should be able to restate all components of the proof of Theorem \ref{cor:diagDimFPAlg} and prove the equality above.

\smallskip
\subsubsection{Noncommutative Cartan pairs and generalised diagonal dimension}

Li, Liao, and Winter showed that finite diagonal dimension of a nondegenerate \cstarSub{} implies that the subalgebra is a \cstarDiag{}.

\begin{thm}[{\cite[Theorem~2.10]{LiLiaoWinter2023diagDim}}] \label{thm:diagDimImpliesCdiagonal}
Let ${(D_A \subseteq A)}$ be a nondegenerate \cstarSub{} with ${\ddim{D_A}{A}<\infty}.$ 
Then $D_A$ is a \cstarDiag{} in $A$.
\end{thm}

An open question is whether a noncommutative version of the above holds. 
More precisely, given a \cstar{} $B$ and
a nondegenerate \cstarSub{} ${(D_A \subseteq A)}$ such that 
    \[\ddimn{D_A}{A}{B}<\infty,\]
does it follow that $D_A$ a noncommutative Cartan subalgebra in $A$?

It is clear that choosing a suitable \cstar{} $B$ as 
coefficients
plays a crucial role in the above question.
This brings us to another direction that can be explored.

\smallskip
\subsubsection{Coefficients in generalised diagonal dimension}

It would be of interest to determine appropriate coefficients $B$ in the definition of generalised diagonal dimension.
Note that in Theorem~\ref{cor:diagDimFPAlg} the \cstar{} $\ell^\infty(\N,\bop{\Hcal})$ is (non\myhyp{}canonically) isomorphic to the \cstarSub{} $\CfpCartan{X}$.
It remains to be seen whether computing the generalised diagonal dimension of a pair ${(D_A \subseteq A)}$ with respect to ${D_A}$, that is, computing ${\ddimn{D_A}{A}{D_A}}$,
yields meaningful results.


%%%%% BIBLIOGRAPHY 
%%%%%%%%%%%%%%%%%%%%%%%%%%%%%%%%%%%%
\bibliographystyle{amsplain}
\bibliography{references.bib}
%%%%%%%%%%%%%%%%%%%%%%%%%%%%%%%%%%%%

\end{document}
